\documentclass[tate,paper=b6j,jafontsize=9pt]{jlreq}

\jlreqsetup{warichu_opening=, warichu_closing=}

\usepackage{ifthen}

\usepackage{mdframed}

\usepackage{graphicx}

\usepackage{lltjp-geometry}

\usepackage[driver=dvipdfmx]{geometry}

\usepackage{pxrubrica}

\usepackage{ulem}

\usepackage{luatexja-fontspec}

\usepackage{newunicodechar}

\usepackage{caption}

\usepackage[kumi=aki, tateaki=1]{kanbun}

%%%%%%%%%%%%%%%%%%%%%%%%%%%%%%
% environments
%%%%%%%%%%%%%%%%%%%%%%%%%%%%%%

% \begin{azconvIndented}{n}{m}はｍ字下げ、改行してｎ字下げ。
% ［＃ｎ字下げ］はｎ字下げ、改行してｎ字下げ。
% ［＃改行天付き］はｎ字下げ、改行して０字下げ。
\newenvironment{azconvIndented}[2]{
  \setlength{\parindent}{#2}
   \leftskip=#1

  }
  {
   \setlength{\parindent}{0\zw}
   \leftskip=0em
  }

% \begin{azconvAlignBottom}{n}は下からｎ字上げ。地付きは下から０字上げ。 
\newenvironment{azconvAlignBottom}[1]{
  \rightskip=#1
 }
 {\rightskip=0\zw}

% #1は字下げ及び字詰めのスコープ内の場合指定される。字下げの場合下げる字数。
% 字詰めの場合は求められる字数に近づくように設定される。
 \newenvironment{azconvFramed}[1][]%
 {\begin{mdframed}[#1,skipabove=8pt]}
{\end{mdframed}}


\newenvironment{azconvNarrow}[1]{
 \rightskip= \dimexpr\linewidth-#1}
{\rightskip=0em}

\newenvironment{azconvFigure}{
 \begin{figure}[h!]
\begin{center}}
{\end{center}
\end{figure}}


%%%%%%%%%%%%%%%%%%%%%%%%%%%%%%
% commands
%%%%%%%%%%%%%%%%%%%%%%%%%%%%%%

%%%%% (re)define various under/overline styles for ulem
%%%%% shamelessly reusing and adjusting code from ulem.sty

\makeatletter

\UL@protected\def\uline{\leavevmode \bgroup
 \UL@setULdepth
 \ifx\UL@on\UL@onin \advance\ULdepth2.8\p@\fi
 \advance\ULdepth-0.5\p@\markoverwith{\lower\ULdepth\hbox
   {\kern-.03em\vbox{\hrule width.2em}\kern-.03em}}%
 \ULon}

\UL@protected\def\oline{\leavevmode \bgroup
 \UL@setULdepth
 \ifx\UL@on\UL@onin \advance\ULdepth2.8\p@\fi
 \advance\ULdepth-0.5\p@\markoverwith{\raise\ULdepth\hbox
   {\kern-.03em\vbox{\hrule width.2em}\kern-.03em}}%
 \ULon}

 \UL@protected\def\uuline{\leavevmode \bgroup
 \UL@setULdepth
 \ifx\UL@on\UL@onin \advance\ULdepthi2.8\p@\fi
 \advance\ULdepth1\p@\markoverwith{\lower\ULdepth\hbox
   {\kern-.03em\vbox{\hrule width.2em\kern1\p@\hrule}\kern-.03em}}%
 \ULon}

\UL@protected\def\ooline{\leavevmode \bgroup
 \UL@setULdepth
 \ifx\UL@on\UL@onin \advance\ULdepth2.8\p@\fi
 \advance\ULdepth-0.5\p@\markoverwith{\raise\ULdepth\hbox
   {\kern-.03em\vbox{\hrule width.2em\kern1\p@\hrule}\kern-.03em}}%
 \ULon}

\UL@protected\def\ooline{\leavevmode \bgroup
 \UL@setULdepth
 \ifx\UL@on\UL@onin \advance\ULdepth2.8\p@\fi
 \advance\ULdepth-.8\p@\markoverwith{\raise\ULdepth\hbox
   {\kern-.03em\vbox{\hrule width.2em\kern1\p@\hrule}\kern-.03em}}%
 \ULon}

\UL@protected\def\dotuline{\leavevmode \bgroup 
  \UL@setULdepth
  \ifx\UL@on\UL@onin \advance\ULdepth2.8\p@\fi
  \markoverwith{\begingroup
   \advance\ULdepth-2\p@\lower\ULdepth\hbox{\kern.06em .\kern.04em}%
     \endgroup}%
  \ULon}

\UL@protected\def\dotoline{\leavevmode \bgroup \UL@setULdepth
 \ifx\UL@on\UL@onin \advance\ULdepth2.8\p@\fi
 \markoverwith{\begingroup
 \advance\ULdepth3\p@\raise\ULdepth\hbox{\kern.06em .\kern.04em}%
	\endgroup}%
 \ULon}

\UL@protected\def\dashuline{\leavevmode \bgroup 
  \UL@setULdepth
  \ifx\UL@on\UL@onin \advance\ULdepth2\p@\fi
  \markoverwith{\begingroup
     \advance\ULdepth-0.8\p@\ 
     \lower\ULdepth\hbox{\kern.06em -\kern.04em}%
     \endgroup}%
  \ULon}

\UL@protected\def\dasholine{\leavevmode \bgroup \UL@setULdepth
 \ifx\UL@on\UL@onin \advance\ULdepth2\p@\fi
 \markoverwith{\begingroup
	\advance\ULdepth1.5\p@\raise\ULdepth\hbox{\kern.06em -\kern.04em}%
	\endgroup}%
  \ULon}

\UL@protected\def\owave{\leavevmode \bgroup 
  \ifdim \ULdepth=\maxdimen \ULdepth 3.5\p@
  \else \advance\ULdepth2\p@ 
  \fi \advance\ULdepth3.2\p@\markoverwith{\raise\ULdepth\hbox{\sixly \char58}}\ULon}
\font\sixly=lasy6 % does not re-load if already loaded, so no memory drain.

\UL@protected\def\uwave{\leavevmode \bgroup 
  \ifdim \ULdepth=\maxdimen \ULdepth 3.5\p@
  \else \advance\ULdepth2\p@ 
  \fi \advance\ULdepth0.5\p@\markoverwith{\lower\ULdepth\hbox{\sixly \char58}}\ULon}
\font\sixly=lasy6 % does not re-load if already loaded, so no memory drain.

\makeatother

%********


\newcommand{\azconvBold}[1]{%
{\textbf{#1}}}

\newcommand{\azconvContributor}[1]{
   \hfill #1

  }

\newcommand{\azconvFbox}[1]{
   \fbox{#1}
  }

\newcommand{\azconvIncludegraphics}[3]{
 \begin{center}
  \hskip -4\zw
  \resizebox{#2}{#1}{\includegraphics[angle=90, origin=c]{#3}}
\end{center}}

\newcommand{\azconvIncludegraphicsInline}[1]{
 \raisebox{-0.6\zw}{\includegraphics[angle=90, origin=c, width=1\zw, height=1\zw]
{#1}}}

\newcommand{\azconvInlineHeader}[2]{%
\hskip 0.5\zw\textbf{#2}\hskip-0.5\zw}

\newcommand{\azconvItalic}[1]{%
{\textit{#1}}}

\newcommand{\azconvKenten}[3][]{%
 \ifthenelse{\equal{#1}{left}}%
 {\kentenmarkintate{#2}\kenten[S]{#3}}%
 {\kentenmarkintate{#2}\kenten{#3}}%
}

\newcommand{\azconvKunten}[1]{%
\raisebox{-0.4\zw}{\scriptsize #1}\normalsize}

%　#1は改ページ、改丁など注記で指示されているもの。デフォルトでは無視する。
\newcommand{\azconvNewpage}[1]{%
 \newpage
}

\newcommand{\azconvOkurigana}[1]{%
\raisebox{0.6\zw}{\scriptsize #1}\normalsize}

\newcommand{\azconvSection}[1]{%
 \section{#1}
}

\newcommand{\azconvSubsection}[1]{%
  \subsection{#1}
 }

\newcommand{\azconvSubsubsection}[1]{%
  \subsubsection{#1}
 }


\newcommand{\azconvSup}[1]{%
\textsuperscript{#1}}

\newcommand{\azconvSub}[1]{%
\textsubscript{#1}}

\newcommand{\azconvSubtitle}[1]{
	{\large
		\begin{azconvIndented}{7\zw}{0\zw}

			#1

	 \end{azconvIndented}
	}
 
  }

\newcommand{\azconvTatechuyoko}[1]{%
 \tatechuyoko{#1}}

\newcommand{\azconvTitle}[1]{
	{\LARGE
	 
	 \begin{azconvIndented}{3\zw}{0\zw}

	 #1
	 
	\end{azconvIndented}
		\vskip 1\zw
	}

  }

 % #1が傍線の種類、#2が左右、#3が対象文字列。
\newcommand{\azconvUline}[3]{%
   \ifthenelse
   {\equal{#1}{solid}}
	{\ifthenelse{
	 \equal{#2}{under}}
	{\oline{#3}}
   {\uline{#3}}}{}%
   %
\ifthenelse
   {\equal{#1}{double}}
	{\ifthenelse{
	 \equal{#2}{under}}
	 {\ooline{#3}}
	{\uuline{#3}}}{}%
	%
\ifthenelse
   {\equal{#1}{dotted}}
	{\ifthenelse{
	 \equal{#2}{under}}
	 {\dotoline{#3}}
	{\dotuline{#3}}}{}%
	%
\ifthenelse
   {\equal{#1}{dashed}}
	{\ifthenelse{
	 \equal{#2}{under}}
	 {\dasholine{#3}}
	{\dashuline{#3}}}{}%
	%
\ifthenelse
   {\equal{#1}{wavy}}
	{\ifthenelse{
	 \equal{#2}{under}}
	 {\owave{#3}}
	{\uwave{#3}}}{}%
}

% #1は必要となるboxのサイズ、#2が文字列。#1はjlreqの場合自動で算出されるので無視。
\newcommand{\azconvWarichu}[2]{
 \warichu{#2}}

% 割中の中で改行が指定されている場合に使う。#1はboxのサイズ、#2が文字列。
\newcommand{\azconvWarichuM}[2]{\linespread{0.5}
{\scriptsize\parbox{#1}{#2}}}

\newcommand{\azconvYokogumi}[1]{%
\raisebox{-.35\zw}{\hbox{\utod #1}}}


 %%%%%%%%%%%%%%%%%%%%%%%%%%%%%%
% for Kanbun kunten 
%
% taken from 漢文マクロ kambunsnx 
%
% 峰山進(閑舎) https://rakunet.org/p/2022/0729
%
%%%%%%%%%%%%%%%%%%%%%%%%%%%%%%

\def\azconvIR{一\kern-.6\zw レ}
\def\azconvNR{二\kern-.4\zw レ}
\def\azconvJR{上\kern-.25\zw レ}
\def\azconvKR{甲\kern-.4\zw レ}


%%%%%%%%%%%%%%%%%%%%%%%%%%%%%%
% setup basic lengths, fonts, etc.
%%%%%%%%%%%%%%%%%%%%%%%%%%%%%%

\setlength{\parindent}{0em}

\renewcommand{\baselinestretch}{1.2}

\captionsetup{font={small, normalfont},labelformat=empty,skip=0pt}

\ModifyHeading{section}{indent=2\zw, format={\Large #2}}
\ModifyHeading{subsection}{indent=3\zw,format={\large #2}}
\ModifyHeading{subsubsection}{indent=4\zw,format={\bfseries #2}}

%\usepackage[noto-jp]{luatexja-preset}

%%%%%%%%%%%%%%%%%%%%%%%%%%%%
%
% Fallback for unavailable characters
%
%%%%%%%%%%%%%%%%%%%%%%%%%%%%


%\newjfontfamily{\fallback}{MingLiU-ExtB}

%\newunicodechar{𥥔}{{\fallback 𥥔}} 


\setlength{\parindent}{0em}

\begin{document}

\thispagestyle{empty}
\vspace*{\fill}
\vskip-3\zw
\azconvTitle{南洲手抄言志録}
\azconvSubtitle{03 南洲手抄言志録}
\azconvContributor{秋月 種樹}
\azconvContributor{佐藤 一斎}
\azconvContributor{西郷 隆盛 編}
\azconvContributor{山田 済斎 訳}
\vfill
\newpage

\begin{azconvIndented}{1\zw}{-1\zw}

一　勿\azconvKunten{下}認\azconvKunten{二}游惰\azconvKunten{一}以爲\azconvKunten{中}寛裕\azconvKunten{上}。勿\azconvKunten{下}認\azconvKunten{二}嚴刻\azconvKunten{一}以爲\azconvKunten{中}直諒\azconvKunten{上}。勿\azconvKunten{下}認\azconvKunten{二}私欲\azconvKunten{一}以爲\azconvKunten{中}志願\azconvKunten{上}。

\end{azconvIndented}
\begin{azconvIndented}{2\zw}{-1\zw}

〔譯〕\jruby[g]{游惰}{いうだ}を\jruby[g]{認}{みと}めて以て\jruby[g]{寛裕}{かんゆう}と爲すこと\jruby[g]{勿}{なか}れ。\jruby[g]{嚴刻}{げんこく}を認めて以て\jruby[g]{直諒}{ちよくりやう}と爲すこと勿れ。\jruby[g]{私欲}{しよく}を認めて以て\jruby[g]{志願}{しぐわん}と爲すこと勿れ。

\end{azconvIndented}

\begin{azconvIndented}{1\zw}{-1\zw}

二　毀譽得喪、眞是人生之雲霧、使\azconvKunten{二}人昏迷\azconvKunten{一}。一\azconvKunten{二}掃此雲霧\azconvKunten{一}、則天青日白。

\end{azconvIndented}
\begin{azconvIndented}{2\zw}{-1\zw}

〔譯〕\jruby[g]{毀譽}{きよ}\jruby[g]{得喪}{とくさう}は、\jruby[g]{眞}{しん}に是れ人生の\jruby[g]{雲霧}{うんむ}、人をして\jruby[g]{昏迷}{こんめい}せしむ。此の雲霧を一\jruby[g]{掃}{さう}せば、則ち\jruby[g]{天}{てん}\jruby[g]{青}{あを}く\jruby[g]{日}{ひ}\jruby[g]{白}{しろ}し。

\end{azconvIndented}
\begin{azconvIndented}{3\zw}{-1\zw}

評徳川\jruby[g]{慶喜}{よしのぶ}公は\jruby[g]{勤王}{きんわう}の臣たり。\jruby[g]{幕吏}{ばくり}の要する所となりて\jruby[g]{朝敵}{てうてき}となる。猶南洲勤王の臣として終りを\jruby[g]{克}{よ}くせざるごとし。公は\jruby[g]{罪}{つみ}を\jruby[g]{宥}{ゆる}し位に\jruby[g]{敍}{じよ}せらる、南洲は永く\jruby[g]{反賊}{はんぞく}の名を\jruby[g]{蒙}{かうむ}る、悲しいかな。（原漢文、下同）

\end{azconvIndented}

\begin{azconvIndented}{1\zw}{-1\zw}

三　唐虞之治、只是情一字。極而言\azconvKunten{レ}之、萬物一體、不\azconvKunten{レ}外\azconvKunten{二}於情之推\azconvKunten{一}。

\end{azconvIndented}
\begin{azconvIndented}{2\zw}{-1\zw}

〔譯〕\jruby[g]{唐虞}{たうぐ}の\jruby[g]{治}{ち}は只是れ情の一字なり。極めて之を言へば、萬物一體も情の\jruby[g]{推}{すゐ}に外ならず。

\end{azconvIndented}
\begin{azconvIndented}{3\zw}{-1\zw}

評南洲、官軍を帥ゐて京師を發す。\jruby[g]{婢}{ひ}あり別れを惜みて\jruby[g]{伏水}{ふしみ}に至る。兵士\jruby[g]{環}{めぐ}つて之を\jruby[g]{視}{み}る。南洲輿中より之を招き、其背を\jruby[g]{拊}{う}つて曰ふ、\jruby[g]{好在}{たつしや}なれと、金を\jruby[g]{懷中}{くわいちゆう}より出して之に與へ、\jruby[g]{旁}{かたは}ら人なき若し。兵士\jruby[g]{太}{はなは}だ其の情を\jruby[g]{匿}{かく}さざるに服す。幕府\jruby[g]{砲臺}{はうだい}を神奈川に\jruby[g]{築}{きづ}き、外人の來り觀るを許さず、木戸公\jruby[g]{役徒}{えきと}に雜り、自ら\jruby[g]{畚}{ふご}を\jruby[g]{荷}{にな}うて之を觀る。茶店の\jruby[g]{老嫗}{らうをう}あり、公の常人に非ざるを知り、善く之を遇す。公志を得るに及んで、厚く之に報ゆ。皆情の\jruby[g]{推}{すゐ}なり。

\end{azconvIndented}

\begin{azconvIndented}{1\zw}{-1\zw}

四　凡作\azconvKunten{レ}事、須\azconvKunten{レ}要\azconvKunten{レ}有\azconvKunten{二}事\azconvKunten{レ}天之心\azconvKunten{一}。不\azconvKunten{レ}要\azconvKunten{レ}有\azconvKunten{二}示\azconvKunten{レ}人之念\azconvKunten{一}。

\end{azconvIndented}
\begin{azconvIndented}{2\zw}{-1\zw}

譯凡そ事を\jruby[g]{作}{な}すには、\jruby[g]{須}{すべか}らく天に\jruby[g]{事}{つか}ふるの心あるを\jruby[g]{要}{えう}すべし。人に示すの\jruby[g]{念}{ねん}あるを要せず。

\end{azconvIndented}

\begin{azconvIndented}{1\zw}{-1\zw}

五　憤一字、是進學機關。舜何人也、予何人也、方是憤。

\end{azconvIndented}
\begin{azconvIndented}{2\zw}{-1\zw}

〔譯〕\jruby[g]{憤}{ふん}の一字、是れ\jruby[g]{進學}{しんがく}の\jruby[g]{機關}{きくわん}なり。\jruby[g]{舜}{しゆん}\jruby[g]{何人}{なんぴと}ぞや、\jruby[g]{予}{われ}何人ぞや、\jruby[g]{方}{まさ}に是れ\jruby[g]{憤}{ふん}。

\end{azconvIndented}

\begin{azconvIndented}{1\zw}{-1\zw}

六　著\azconvKunten{レ}眼高、則見\azconvKunten{レ}理不\azconvKunten{レ}岐。

\end{azconvIndented}
\begin{azconvIndented}{2\zw}{-1\zw}

〔譯〕\jruby[g]{眼}{がん}を\jruby[g]{著}{つ}くること高ければ、則ち\jruby[g]{理}{り}を見ること\jruby[g]{岐}{き}せず。

\end{azconvIndented}
\begin{azconvIndented}{3\zw}{-1\zw}

評三條公は西三條、東久世諸公と長門に走る、之を七\jruby[g]{卿}{きやう}\jruby[g]{脱走}{だつさう}と謂ふ。幕府之を\jruby[g]{宰府}{ざいふ}に\jruby[g]{竄}{ざん}す。既にして七卿が勤王の士を\jruby[g]{募}{つの}り國家を亂さんと欲するを憂へ、\jruby[g]{浪華}{なには}に\jruby[g]{幽}{いう}するの\jruby[g]{議}{ぎ}あり。南洲等\jruby[g]{力}{つと}めて之を拒ぎ、事終に\jruby[g]{熄}{や}む。南洲人に\jruby[g]{語}{かた}つて曰ふ、七卿中他日\jruby[g]{關白}{くわんぱく}に任ぜらるゝ者は、必三條公ならんと、果して然りき。

\end{azconvIndented}

\begin{azconvIndented}{1\zw}{-1\zw}

七　性同而質異。質異、教之所\azconvKunten{二}由設\azconvKunten{一}也。性同、教之所\azconvKunten{二}由立\azconvKunten{一}也。

\end{azconvIndented}
\begin{azconvIndented}{2\zw}{-1\zw}

〔譯〕\jruby[g]{性}{せい}は同じうして而て\jruby[g]{質}{しつ}は\jruby[g]{異}{ことな}る。質異るは\jruby[g]{教}{をしへ}の由つて\jruby[g]{設}{まう}けらるゝ所なり。性同じきは教の由つて立つ所なり。

\end{azconvIndented}

\begin{azconvIndented}{1\zw}{-1\zw}

八　喪\azconvKunten{レ}己斯喪\azconvKunten{レ}人。喪\azconvKunten{レ}人斯喪\azconvKunten{レ}物。

\end{azconvIndented}
\begin{azconvIndented}{2\zw}{-1\zw}

〔譯〕\jruby[g]{己}{おのれ}を\jruby[g]{喪}{うしな}へば\jruby[g]{斯}{こゝ}に人を\jruby[g]{喪}{うしな}ふ。人を喪へば斯に\jruby[g]{物}{もの}を喪ふ。

\end{azconvIndented}

\begin{azconvIndented}{1\zw}{-1\zw}

九　士貴\azconvKunten{二}獨立自信\azconvKunten{一}矣。依\azconvKunten{レ}熱附\azconvKunten{レ}炎之念、不\azconvKunten{レ}可\azconvKunten{レ}起。

\end{azconvIndented}
\begin{azconvIndented}{2\zw}{-1\zw}

〔譯〕\jruby[g]{士}{し}は\jruby[g]{獨立}{どくりつ}\jruby[g]{自信}{じしん}を\jruby[g]{貴}{たふと}ぶ。\jruby[g]{熱}{ねつ}に\jruby[g]{依}{よ}り\jruby[g]{炎}{えん}に\jruby[g]{附}{つ}くの\jruby[g]{念}{ねん}、起す可らず。

\end{azconvIndented}
\begin{azconvIndented}{3\zw}{-1\zw}

〔評〕\jruby[g]{慶應}{けいおう}三年九月、山内\jruby[g]{容堂}{ようだう}公は寺村\jruby[g]{左膳}{さぜん}、後藤\jruby[g]{象}{しやう}次郎を以て使となし、書を幕府に\jruby[g]{呈}{てい}す。曰ふ、中古以\jruby[g]{還}{くわん}、\jruby[g]{政刑}{せいけい}武門に出づ。洋人來航するに及んで、\jruby[g]{物議}{ぶつぎ}\jruby[g]{紛々}{ふん{〳〵}}、東攻西\jruby[g]{撃}{げき}して、\jruby[g]{内訌}{ないこう}嘗て\jruby[g]{戢}{をさま}る時なく、終に外國の\jruby[g]{輕侮}{けいぶ}を\jruby[g]{招}{まね}くに至る。此れ\jruby[g]{政令}{せいれい}二\jruby[g]{途}{と}に出で、天下耳目の\jruby[g]{屬}{ぞく}する所を異にするが故なり。今や時勢一\jruby[g]{變}{ぺん}して\jruby[g]{舊規}{きうき}を\jruby[g]{墨守}{ぼくしゆ}す可らず、宜しく政\jruby[g]{權}{けん}を王室に還し、以て萬國\jruby[g]{竝立}{へいりつ}の\jruby[g]{基礎}{きそ}を建つべし。其れ則ち當今の\jruby[g]{急務}{きふむ}にして、而て容堂の\jruby[g]{至願}{しぐわん}なり。\jruby[g]{幕}{ばく}下の\jruby[g]{賢}{けん}なる、必之を\jruby[g]{察}{さつ}するあらんと。他日幕府の政權を\jruby[g]{還}{かへ}せる、其事實に公の\jruby[g]{呈書}{ていしよ}に\jruby[g]{本}{もと}づけり。當時\jruby[g]{幕府}{ばくふ}既に\jruby[g]{衰}{おとろ}へたりと雖、\jruby[g]{威權}{ゐけん}未だ地に\jruby[g]{墜}{お}ちず。公\jruby[g]{抗論}{かうろん}して\jruby[g]{忌}{い}まず、獨立の見ありと謂ふべし。

\end{azconvIndented}

\begin{azconvIndented}{1\zw}{-1\zw}

一〇　有\azconvKunten{二}本然之眞己\azconvKunten{一}、有\azconvKunten{二}躯殼之假己\azconvKunten{一}。須\azconvKunten{レ}要\azconvKunten{二}自認得\azconvKunten{一}。

\end{azconvIndented}
\begin{azconvIndented}{2\zw}{-1\zw}

〔譯〕\jruby[g]{本然}{ほんぜん}の\jruby[g]{眞己}{しんこ}有り、\jruby[g]{躯殼}{くかく}の\jruby[g]{假己}{かこ}有り。須らく自ら\jruby[g]{認}{みと}め得んことを要すべし。

\end{azconvIndented}
\begin{azconvIndented}{3\zw}{-1\zw}

評南洲\jruby[g]{胃}{い}を病む。英醫\jruby[g]{偉利斯}{いりす}之を\jruby[g]{診}{しん}して、\jruby[g]{勞動}{らうどう}を\jruby[g]{勸}{すゝ}む。南洲是より山野に\jruby[g]{游獵}{いうれふ}せり。人或は病なくして犬を\jruby[g]{牽}{ひ}き兎を\jruby[g]{逐}{お}ひ、自ら南洲を學ぶと謂ふ、\jruby[g]{疎}{そ}なり。

\end{azconvIndented}

\begin{azconvIndented}{1\zw}{-1\zw}

一一　雲煙聚\azconvKunten{二}於不\azconvKunten{\azconvIR}得\azconvKunten{レ}已。風雨洩\azconvKunten{二}於不\azconvKunten{\azconvIR}得\azconvKunten{レ}已。雷霆震\azconvKunten{二}於不\azconvKunten{\azconvIR}得\azconvKunten{レ}已。斯可\azconvKunten{三}以觀\azconvKunten{二}至誠之作用\azconvKunten{一}。

\end{azconvIndented}
\begin{azconvIndented}{2\zw}{-1\zw}

〔譯〕\jruby[g]{雲煙}{うんえん}は\jruby[g]{已}{や}むことを得ざるに\jruby[g]{聚}{あつま}る。\jruby[g]{風雨}{ふうう}は已むことを得ざるに\jruby[g]{洩}{も}る。\jruby[g]{雷霆}{らいてい}は已むことを得ざるに\jruby[g]{震}{ふる}ふ。\jruby[g]{斯}{こゝ}に以て\jruby[g]{至誠}{しせい}の\jruby[g]{作用}{さよう}を\jruby[g]{觀}{み}る可し。

\end{azconvIndented}

\begin{azconvIndented}{1\zw}{-1\zw}

一二　動\azconvKunten{二}於不\azconvKunten{レ}得\azconvKunten{レ}已之勢\azconvKunten{一}、則動而不\azconvKunten{レ}括。履\azconvKunten{二}於不\azconvKunten{レ}可\azconvKunten{レ}枉之途\azconvKunten{一}、則履而不\azconvKunten{レ}危。

\end{azconvIndented}
\begin{azconvIndented}{2\zw}{-1\zw}

譯已むことを得ざるの\jruby[g]{勢}{いきほひ}に\jruby[g]{動}{うご}けば、則ち動いて\jruby[g]{括}{くわつ}せず。\jruby[g]{枉}{ま}ぐ可らざるの\jruby[g]{途}{みち}を\jruby[g]{履}{ふ}めば、則ち履んで\jruby[g]{危}{あやふ}からず。

\end{azconvIndented}
\begin{azconvIndented}{3\zw}{-1\zw}

評官軍江戸を\jruby[g]{伐}{う}つ、關西諸侯兵を出して之に從ふ。是より先き\jruby[g]{尾藩}{びはん}\jruby[g]{宗家}{そうけ}を\jruby[g]{援}{たす}けんと欲する者ありて、\jruby[g]{私}{ひそ}かに\jruby[g]{聲息}{せいそく}を江戸に\jruby[g]{通}{つう}ず。\jruby[g]{尾}{び}公之を\jruby[g]{患}{うれ}へ、田中\jruby[g]{不二麿}{ふじまろ}、丹羽淳太郎等と議して、大義\jruby[g]{親}{しん}を\jruby[g]{滅}{ほろぼ}すの令を下す、實に已むことを得ざるの\jruby[g]{擧}{きよ}に出づ。一藩の\jruby[g]{方向}{はうかう}以て定れり。

\end{azconvIndented}

\begin{azconvIndented}{1\zw}{-1\zw}

一三　聖人如\azconvKunten{二}強健無\azconvKunten{レ}病人\azconvKunten{一}。賢人如\azconvKunten{二}攝生愼\azconvKunten{レ}病人\azconvKunten{一}。常人如\azconvKunten{二}虚羸多\azconvKunten{レ}病人\azconvKunten{一}。

\end{azconvIndented}
\begin{azconvIndented}{2\zw}{-1\zw}

譯聖人は\jruby[g]{強健}{きやうけん}病無き人の如し。賢人は\jruby[g]{攝生}{せつしやう}病を\jruby[g]{愼}{つゝし}む人の如し。常人は\jruby[g]{虚羸}{きよるゐ}病多き人の如し。

\end{azconvIndented}

\begin{azconvIndented}{1\zw}{-1\zw}

一四　急迫敗\azconvKunten{レ}事。寧耐成\azconvKunten{レ}事。

\end{azconvIndented}
\begin{azconvIndented}{2\zw}{-1\zw}

〔譯〕\jruby[g]{急迫}{きふはく}は事を\jruby[g]{敗}{やぶ}る。\jruby[g]{寧耐}{ねいたい}は事を\jruby[g]{成}{な}す。

\end{azconvIndented}
\begin{azconvIndented}{3\zw}{-1\zw}

評大坂城\jruby[g]{陷}{おちい}る。徳川\jruby[g]{慶喜}{よしのぶ}公火船に乘りて江戸に歸り、諸侯を召して罪を\jruby[g]{俟}{ま}つの状を告ぐ。余時に江戸に在り、特に\jruby[g]{別廳}{べつちやう}に\jruby[g]{召}{め}し告げて曰ふ。事此に至る、言ふ可きなし。汝將に京に入らんとすと\jruby[g]{聞}{き}く、請ふ吾が爲めに\jruby[g]{恭順}{きようじゆん}の意を致せと。余江戸を發して桑名に\jruby[g]{抵}{いた}り、柳原\jruby[g]{前光}{さきみつ}公軍を\jruby[g]{督}{とく}して至るに遇ふ。余爲めに之を告ぐ。京師に至るに及んで、松平\jruby[g]{春嶽}{しゆんがく}公を見て又之を告ぐ。慶喜公江戸城に在り、衆皆之に\jruby[g]{逼}{せま}り、死を以て城を守らんことを請ふ。公\jruby[g]{聽}{き}かず、水戸に赴く、近臣二三十名從ふ。衆奉じて以て主と爲すべきものなく、或は\jruby[g]{散}{さん}じて四方に\jruby[g]{之}{ゆ}き、或は\jruby[g]{上野}{うへの}に\jruby[g]{據}{よ}る。若し公をして\jruby[g]{耐忍}{たいにん}の力無く、共に\jruby[g]{怒}{いか}つて事を擧げしめば、則ち府下悉く\jruby[g]{焦土}{せうど}と爲らん。\jruby[g]{假令}{たとひ}都を遷すも、其の盛大を\jruby[g]{極}{きは}むること今日の如きは實に難からん。然らば則ち公常人の\jruby[g]{忍}{しの}ぶ能はざる所を忍ぶ、其功亦多し。\jruby[g]{舊}{きう}藩士\jruby[g]{日高誠實}{ひだかせいじつ}時に句あり云ふ。

\end{azconvIndented}
\begin{azconvIndented}{3\zw}{0\zw}

「\jruby[g]{功烈}{こうれつ}尤も多かりしは\jruby[g]{前内府}{ぜんないふ}。\jruby[g]{至尊}{しそん}直に\jruby[g]{鶴城}{かくじやう}の中に在り」と。

\end{azconvIndented}

\begin{azconvIndented}{1\zw}{-1\zw}

一五　聖人安\azconvKunten{レ}死。賢人分\azconvKunten{レ}死。常人恐\azconvKunten{レ}死。

\end{azconvIndented}
\begin{azconvIndented}{2\zw}{-1\zw}

譯聖人は死を\jruby[g]{安}{やす}んず。賢人は死を\jruby[g]{分}{ぶん}とす。常人は死を\jruby[g]{恐}{おそ}る。

\end{azconvIndented}

\begin{azconvIndented}{1\zw}{-1\zw}

一六　賢者臨\azconvKunten{レ}※{\small ［＃※は「歹＋勿」、33-1］}、見\azconvKunten{二}理當\azconvKunten{\azconvIR}然、以爲\azconvKunten{レ}分、恥\azconvKunten{レ}畏\azconvKunten{レ}死、而希\azconvKunten{レ}安\azconvKunten{レ}死、故神氣不\azconvKunten{レ}亂。又有\azconvKunten{二}遺訓\azconvKunten{一}、足\azconvKunten{二}以聳\azconvKunten{\azconvIR}聽。而其不\azconvKunten{レ}及\azconvKunten{二}聖人\azconvKunten{一}亦在\azconvKunten{二}於此\azconvKunten{一}。聖人平生言動無\azconvKunten{二}一非\azconvKunten{\azconvIR}訓。而臨\azconvKunten{レ}※{\small ［＃※は「歹＋勿」、33-3］}、未\azconvKunten{三}必爲\azconvKunten{二}遺訓\azconvKunten{一}。視\azconvKunten{二}死生\azconvKunten{一}眞如\azconvKunten{二}晝夜\azconvKunten{一}、無\azconvKunten{レ}所\azconvKunten{レ}著\azconvKunten{レ}念。

\end{azconvIndented}
\begin{azconvIndented}{2\zw}{-1\zw}

譯賢者は\jruby[g]{{※}}{ぼつ}{\small ［＃※は「歹＋勿」、33-4］}するに\jruby[g]{臨}{のぞ}み、\jruby[g]{理}{り}の\jruby[g]{當}{まさ}に然るべきを見て、以て\jruby[g]{分}{ぶん}と爲し、死を\jruby[g]{畏}{おそ}るゝを\jruby[g]{恥}{は}ぢて、死を\jruby[g]{安}{やす}んずるを\jruby[g]{希}{こひねが}ふ、故に\jruby[g]{神氣}{しんき}\jruby[g]{亂}{みだ}れず。又\jruby[g]{遺訓}{いくん}あり、以て\jruby[g]{聽}{ちやう}を\jruby[g]{聳}{そびや}かすに足る。而かも其の聖人に及ばざるも亦此に在り。聖人は平生の\jruby[g]{言動}{げんどう}一として訓に非ざるは無し。而て※{\small ［＃※は「歹＋勿」、33-6］}するに\jruby[g]{臨}{のぞ}みて、未だ必しも遺訓を\jruby[g]{爲}{つく}らず。\jruby[g]{死生}{しせい}を\jruby[g]{視}{み}ること眞に\jruby[g]{晝夜}{ちうや}の如し、\jruby[g]{念}{ねん}を\jruby[g]{著}{つ}くる所無し。

\end{azconvIndented}
\begin{azconvIndented}{3\zw}{-1\zw}

評十年の\jruby[g]{役}{えき}、私學校の\jruby[g]{徒}{と}、\jruby[g]{彈藥}{だんやくせ}\jruby[g]{製造所}{いざうじよ}を\jruby[g]{掠}{かす}む。南洲時に兎を\jruby[g]{大隈}{おほすみ}山中に\jruby[g]{逐}{お}ふ。之を聞いて\jruby[g]{猝}{にはか}に\jruby[g]{色}{いろ}を\jruby[g]{變}{か}へて曰ふ、\jruby[g]{誤}{しま}つたと。\jruby[g]{爾後}{じご}肥後日向に轉戰して、神色\jruby[g]{夷然}{いぜん}たり。

\end{azconvIndented}

\begin{azconvIndented}{1\zw}{-1\zw}

一七　堯舜文王、其所\azconvKunten{レ}遺典謨訓誥、皆可\azconvKunten{三}以爲\azconvKunten{二}萬世法\azconvKunten{一}。何遺命如\azconvKunten{レ}之。至\azconvKunten{二}於成王顧命、曾子善言\azconvKunten{一}、賢人分上自當\azconvKunten{レ}如\azconvKunten{レ}此已。因疑孔子泰山之歌、後人假託爲\azconvKunten{レ}之。檀弓叵\azconvKunten{レ}信、多\azconvKunten{二}此類\azconvKunten{一}。欲\azconvKunten{レ}尊\azconvKunten{二}聖人\azconvKunten{一}、而却爲\azconvKunten{二}之累\azconvKunten{一}。

\end{azconvIndented}
\begin{azconvIndented}{2\zw}{-1\zw}

〔譯〕\jruby[g]{堯舜}{げうしゆん}文王は、其の\jruby[g]{遺}{のこ}す所の\jruby[g]{典謨}{てんぼ}\jruby[g]{訓誥}{くんかう}、皆以て萬世の法と爲す可し。何の\jruby[g]{遺命}{いめい}か之に\jruby[g]{如}{し}かん。\jruby[g]{成}{せい}王の\jruby[g]{顧命}{こめい}、\jruby[g]{曾}{そう}子の善言に至つては、賢人の\jruby[g]{分}{ぶん}上\jruby[g]{自}{おのづか}ら\jruby[g]{當}{まさ}に此の如くなるべきのみ。因つて\jruby[g]{疑}{うたが}ふ、孔子\jruby[g]{泰山}{たいざん}の歌、後人\jruby[g]{假託}{かたく}之を\jruby[g]{爲}{つく}れるならん。\jruby[g]{檀弓}{だんぐう}の信じ\jruby[g]{叵}{がた}きこと此の類多し。聖人を尊ばんと欲して、\jruby[g]{却}{かへ}つて之が\jruby[g]{累}{るゐ}を爲せり。

\end{azconvIndented}

\begin{azconvIndented}{1\zw}{-1\zw}

一八　一部歴史、皆傳\azconvKunten{二}形迹\azconvKunten{一}、而情實或不\azconvKunten{レ}傳。讀\azconvKunten{レ}史者、須\azconvKunten{レ}要\azconvKunten{下}就\azconvKunten{二}形迹\azconvKunten{一}以討\azconvKunten{中}出情實\azconvKunten{上}。

\end{azconvIndented}
\begin{azconvIndented}{2\zw}{-1\zw}

譯一\jruby[g]{部}{ぶ}の\jruby[g]{歴史}{れきし}、皆\jruby[g]{形迹}{けいせき}を\jruby[g]{傳}{つた}へて、\jruby[g]{情實}{じやうじつ}或は傳らず。史を讀む者は、須らく形迹に\jruby[g]{就}{つ}いて以て情實を\jruby[g]{討}{たづ}ね出だすことを要すべし。

\end{azconvIndented}

\begin{azconvIndented}{1\zw}{-1\zw}

一九　博聞強記、聰明横也。精義入\azconvKunten{レ}神、聰明竪也。

\end{azconvIndented}
\begin{azconvIndented}{2\zw}{-1\zw}

〔譯〕\jruby[g]{博聞強記}{はくぶんきやうき}は、\jruby[g]{聰明}{そうめい}の\jruby[g]{横}{よこ}なり。\jruby[g]{精義}{せいぎ}神に入るは、\jruby[g]{聰明}{そうめい}の\jruby[g]{竪}{たて}なり。

\end{azconvIndented}

\begin{azconvIndented}{1\zw}{-1\zw}

二〇　生物皆畏\azconvKunten{レ}死。人其靈也、當\azconvKunten{下}從\azconvKunten{二}畏\azconvKunten{レ}死之中\azconvKunten{一}、揀\azconvKunten{中}出不\azconvKunten{レ}畏\azconvKunten{レ}死之理\azconvKunten{上}。吾思、我身天物也。死生之權在\azconvKunten{レ}天、當\azconvKunten{レ}順\azconvKunten{二}受之\azconvKunten{一}。我之生也、自然而生、生時未\azconvKunten{二}嘗知\azconvKunten{\azconvIR}喜矣。則我之死也、應\azconvKunten{三}亦自然而死、死時未\azconvKunten{二}嘗知\azconvKunten{\azconvIR}悲也。天生\azconvKunten{レ}之而天死\azconvKunten{レ}之、一聽\azconvKunten{二}于天\azconvKunten{一}而已、吾何畏焉。吾性即天也。躯殼則藏\azconvKunten{レ}天之室也。精氣之爲\azconvKunten{レ}物也、天寓\azconvKunten{二}於此室\azconvKunten{一}。遊魂之爲\azconvKunten{レ}變也、天離\azconvKunten{二}於此室\azconvKunten{一}。死之後即生之前、生之前即死之後。而吾性之所\azconvKunten{二}以爲\azconvKunten{\azconvIR}性者、恒在\azconvKunten{二}於死生之外\azconvKunten{一}、吾何畏焉。夫晝夜一理、幽明一理。原\azconvKunten{レ}始反\azconvKunten{レ}終、知\azconvKunten{二}死生之理\azconvKunten{一}、何其易簡而明白也。吾人當\azconvKunten{下}以\azconvKunten{二}此理\azconvKunten{一}自省\azconvKunten{上}焉。

\end{azconvIndented}
\begin{azconvIndented}{2\zw}{-1\zw}

譯生物は皆死を\jruby[g]{畏}{おそ}る。人は其\jruby[g]{靈}{れい}なり、當に死を畏るゝの中より死を畏れざるの理を\jruby[g]{揀出}{けんしゆつ}すべし。吾れ思ふ、我が身は天物なり。死生の\jruby[g]{權}{けん}は天に在り、當に之を\jruby[g]{順受}{じゆんじゆ}すべし。我れの生るゝや自然にして生る、生るゝ時未だ嘗て\jruby[g]{喜}{よろこ}ぶことを知らず。則ち我の死するや\jruby[g]{應}{まさ}に亦自然にして死し、死する時未だ嘗て悲むことを知らざるべし。天之を生みて、天之を\jruby[g]{死}{ころ}す、一に天に\jruby[g]{聽}{まか}さんのみ、吾れ何ぞ畏れん。吾が性は即ち天なり、\jruby[g]{躯殼}{くかく}は則ち天を\jruby[g]{藏}{おさ}むるの室なり。\jruby[g]{精氣}{せいき}の物と爲るや、天此の室に\jruby[g]{寓}{ぐう}す。\jruby[g]{遊魂}{いうこん}の\jruby[g]{變}{へん}を爲すや、天此の室を\jruby[g]{離}{はな}る。死の後は即ち生の前なり、生の前は即ち死の後なり。而て吾が性の性たる所以は、\jruby[g]{恒}{つね}に死生の外に在り、吾れ何ぞ畏れん。夫れ晝夜は一\jruby[g]{理}{り}なり、\jruby[g]{幽明}{いうめい}は一理なり。始めを\jruby[g]{原}{たづ}ねて\jruby[g]{終}{をは}りに\jruby[g]{反}{かへ}らば、死生の理を知る、何ぞ其の\jruby[g]{易簡}{いかん}にして明白なるや。吾人は當に此の理を以て\jruby[g]{自省}{じせう}すべし。

\end{azconvIndented}

\begin{azconvIndented}{1\zw}{-1\zw}

二一　畏\azconvKunten{レ}死者生後之情也、有\azconvKunten{二}躯殼\azconvKunten{一}而後有\azconvKunten{二}是情\azconvKunten{一}。不\azconvKunten{レ}畏\azconvKunten{レ}死者生前之性也、離\azconvKunten{二}躯殼\azconvKunten{一}而始見\azconvKunten{二}是性\azconvKunten{一}。人須\azconvKunten{レ}自\azconvKunten{二}得不\azconvKunten{レ}畏\azconvKunten{レ}死之理於畏\azconvKunten{レ}死之中\azconvKunten{一}、庶\azconvKunten{二}乎復\azconvKunten{\azconvIR}性焉。

\end{azconvIndented}
\begin{azconvIndented}{2\zw}{-1\zw}

譯死を畏るゝは生後の情なり、\jruby[g]{躯殼}{くかく}有つて後に\jruby[g]{是}{こ}の情あり。死を畏れざるは生前の性なり、\jruby[g]{躯殼}{くかく}を\jruby[g]{離}{はな}れて始て是の性を見る。人は\jruby[g]{須}{すべか}らく死を畏れざるの理を死を畏るゝの中に\jruby[g]{自得}{じとく}すべし、性に\jruby[g]{復}{かへ}るに\jruby[g]{庶}{ちか}し。

\end{azconvIndented}
\begin{azconvIndented}{3\zw}{-1\zw}

評幕府勤王の士を\jruby[g]{逮}{とら}ふ。南洲及び\jruby[g]{伊地知正治}{いぢちまさはる}、\jruby[g]{海江田武治}{かいえだたけはる}等尤も其の\jruby[g]{指目}{しもく}する所となる。僧\jruby[g]{月照}{げつせう}嘗て近衞公の\jruby[g]{密命}{みつめい}を\jruby[g]{喞}{ふく}みて水戸に至る、幕吏之を\jruby[g]{索}{もと}むること急なり。南洲其の免れざることを知り相共に鹿兒島に\jruby[g]{奔}{はし}る。一日南洲、月照の宅を\jruby[g]{訪}{と}ふ。此の夜月色\jruby[g]{清輝}{せいき}なり。\jruby[g]{預}{あらかじ}め\jruby[g]{酒饌}{しゆせん}を\jruby[g]{具}{そな}へ、舟を薩海に\jruby[g]{泛}{うか}ぶ、南洲及び平野次郎一僕と從ふ。月照船頭に立ち、和歌を朗吟して南洲に示す、南洲\jruby[g]{首肯}{しゆかう}する所あるものゝ如し、遂に相\jruby[g]{擁}{よう}して海に\jruby[g]{投}{とう}ず。次郎等水聲起るを聞いて、\jruby[g]{倉皇}{さうくわう}として之を救ふ。月照既に死して、南洲は\jruby[g]{蘇}{よみがへ}ることを得たり。南洲は\jruby[g]{終身}{しゆうしん}月照と死せざりしを\jruby[g]{憾}{うら}みたりと云ふ。

\end{azconvIndented}

\begin{azconvIndented}{1\zw}{-1\zw}

二二　誘掖而導\azconvKunten{レ}之、教之常也。警戒而喩\azconvKunten{レ}之、教之時也。躬行以率\azconvKunten{レ}之、教之本也。不\azconvKunten{レ}言而化\azconvKunten{レ}之、教之神也。抑而揚\azconvKunten{レ}之、激而進\azconvKunten{レ}之、教之權而變也。教亦多\azconvKunten{レ}術矣。

\end{azconvIndented}
\begin{azconvIndented}{2\zw}{-1\zw}

〔譯〕\jruby[g]{誘掖}{いうえき}して之を\jruby[g]{導}{みちび}くは、教の常なり。\jruby[g]{警戒}{けいかい}して之を\jruby[g]{喩}{さと}すは、教の時なり。\jruby[g]{躬}{み}に行うて之を\jruby[g]{率}{ひ}きゐるは、教の本なり。言はずして之を化するは、教の\jruby[g]{神}{しん}なり。\jruby[g]{抑}{おさ}へて之を\jruby[g]{揚}{あ}げ、\jruby[g]{激}{げき}して之を\jruby[g]{進}{すゝ}ましむるは、教の\jruby[g]{權}{けん}にして而て\jruby[g]{變}{へん}なり。教も亦\jruby[g]{術}{じゆつ}多し。

\end{azconvIndented}

\begin{azconvIndented}{1\zw}{-1\zw}

二三　閑想客感、由\azconvKunten{二}志之不\azconvKunten{\azconvIR}立。一志既立、百邪退聽。譬\azconvKunten{二}之清泉湧出、旁水不\azconvKunten{\azconvIR}得\azconvKunten{二}渾入\azconvKunten{一}。

\end{azconvIndented}
\begin{azconvIndented}{2\zw}{-1\zw}

〔譯〕\jruby[g]{閑想}{かんさう}\jruby[g]{客感}{きやくかん}は、志の立たざるに由る。一志既に立てば、百邪退き\jruby[g]{聽}{き}く。之を\jruby[g]{清泉}{せいせん}\jruby[g]{湧出}{ようしゆつ}せば、\jruby[g]{旁水}{ばうすゐ}\jruby[g]{渾入}{こんにふ}することを得ざるに\jruby[g]{譬}{たと}ふべし。

\end{azconvIndented}
\begin{azconvIndented}{3\zw}{-1\zw}

評政府\jruby[g]{郡縣}{ぐんけん}の\jruby[g]{治}{ち}を\jruby[g]{復}{ふく}せんと欲す、木戸公と南洲と尤も之を主張す。或ひと南洲を見て之を説く、南洲曰く\jruby[g]{諾}{だく}すと。其人又之を説く、南洲曰く、吉之助の一諾、死以て之を守ると、\jruby[g]{他語}{たご}を\jruby[g]{交}{まじ}へず。

\end{azconvIndented}

\begin{azconvIndented}{1\zw}{-1\zw}

二四　心爲\azconvKunten{レ}靈。其條理動\azconvKunten{二}於情識\azconvKunten{一}、謂\azconvKunten{二}之欲\azconvKunten{一}。欲有\azconvKunten{二}公私\azconvKunten{一}、情識之通\azconvKunten{二}於條理\azconvKunten{一}爲\azconvKunten{レ}公。條理之滯\azconvKunten{二}於情識\azconvKunten{一}爲\azconvKunten{レ}私。自辨\azconvKunten{二}其通滯\azconvKunten{一}者、即便心之靈。

\end{azconvIndented}
\begin{azconvIndented}{2\zw}{-1\zw}

譯心を\jruby[g]{靈}{れい}と爲す。其の\jruby[g]{條理}{でうり}の\jruby[g]{情識}{じやうしき}に\jruby[g]{動}{うご}く、之を\jruby[g]{欲}{よく}と謂ふ。欲に\jruby[g]{公私}{こうし}有り、情識の條理に通ずるを公と爲す。條理の情識に\jruby[g]{滯}{とゞこほ}るを私と爲す。自ら其の\jruby[g]{通}{つう}と\jruby[g]{滯}{たい}とを\jruby[g]{辨}{べん}ずるは、即ち心の\jruby[g]{靈}{れい}なり。

\end{azconvIndented}

\begin{azconvIndented}{1\zw}{-1\zw}

二五　人一生所\azconvKunten{レ}遭、有\azconvKunten{二}險阻\azconvKunten{一}、有\azconvKunten{二}坦夷\azconvKunten{一}、有\azconvKunten{二}安流\azconvKunten{一}、有\azconvKunten{二}驚瀾\azconvKunten{一}。是氣數自然、竟不\azconvKunten{レ}能\azconvKunten{レ}免、即易理也。人宜\azconvKunten{二}居而安、玩而樂\azconvKunten{一}焉。若趨\azconvKunten{二}避之\azconvKunten{一}、非\azconvKunten{二}達者之見\azconvKunten{一}。

\end{azconvIndented}
\begin{azconvIndented}{2\zw}{-1\zw}

譯人一生\jruby[g]{遭}{あ}ふ所、\jruby[g]{險阻}{けんそ}有り、\jruby[g]{坦夷}{たんい}有り、\jruby[g]{安流}{あんりう}有り、\jruby[g]{驚瀾}{きやうらん}有り。是れ\jruby[g]{氣數}{きすう}の自然にして、\jruby[g]{竟}{つひ}に\jruby[g]{免}{まぬが}るゝ能はず、即ち\jruby[g]{易理}{えきり}なり。人宜しく居つて安んじ、\jruby[g]{玩}{もてあそ}んで\jruby[g]{樂}{たの}しむべし。若し之を\jruby[g]{趨避}{すうひ}せば、\jruby[g]{達}{たつ}者の見に非ず。

\end{azconvIndented}
\begin{azconvIndented}{3\zw}{-1\zw}

評或ひと岩倉公幕を佐くと\jruby[g]{讒}{ざん}す。公\jruby[g]{薙髮}{ていはつ}して岩倉邸に\jruby[g]{蟄居}{ちつきよ}す。大橋\jruby[g]{愼藏}{しんざう}、\jruby[g]{香}{か}川\jruby[g]{敬}{けい}三、玉松\jruby[g]{操}{みさを}、北島\jruby[g]{秀朝}{ひでとも}等、公の志を知り、深く\jruby[g]{結納}{けつなふ}す。南洲及び大久保公、木戸公、後藤象次郎、坂本龍馬等公を洛東より迎へて、朝政に任ぜしむ。公既に職に在り、\jruby[g]{屡}{しば{〳〵}}\jruby[g]{刺客}{せきかく}の\jruby[g]{狙撃}{そげき}する所となり、\jruby[g]{危難}{きなん}\jruby[g]{累}{しき}りに至る、而かも\jruby[g]{毫}{がう}も\jruby[g]{趨避}{すうひ}せず。

\end{azconvIndented}

\begin{azconvIndented}{1\zw}{-1\zw}

二六　心之官則思。思字只是工夫字。思則愈精明、愈篤實。自\azconvKunten{二}其篤實\azconvKunten{一}謂\azconvKunten{二}之行\azconvKunten{一}、自\azconvKunten{二}其精明\azconvKunten{一}謂\azconvKunten{二}之知\azconvKunten{一}。知行歸\azconvKunten{二}於一思字\azconvKunten{一}。

\end{azconvIndented}
\begin{azconvIndented}{2\zw}{-1\zw}

譯心の\jruby[g]{官}{かん}は則ち思ふ。思の字只是れ\jruby[g]{工夫}{くふう}の字なり。思へば則ち愈\jruby[g]{精明}{せいめい}なり、愈\jruby[g]{篤實}{とくじつ}なり。其の篤實より之を行と謂ひ、其の精明より之を知と謂ふ。知と行とは一の思の字に\jruby[g]{歸}{き}す。

\end{azconvIndented}

\begin{azconvIndented}{1\zw}{-1\zw}

二七　處\azconvKunten{レ}晦者能見\azconvKunten{レ}顯。據\azconvKunten{レ}顯者不\azconvKunten{レ}見\azconvKunten{レ}晦。

\end{azconvIndented}
\begin{azconvIndented}{2\zw}{-1\zw}

〔譯〕\jruby[g]{晦}{くわい}に\jruby[g]{處}{を}る者は能く\jruby[g]{顯}{けん}を見る。顯に\jruby[g]{據}{よ}る者は晦を見ず。

\end{azconvIndented}

\begin{azconvIndented}{1\zw}{-1\zw}

二八　取\azconvKunten{二}信於人\azconvKunten{一}難也。人不\azconvKunten{レ}信\azconvKunten{二}於口\azconvKunten{一}、而信\azconvKunten{二}於躬\azconvKunten{一}。不\azconvKunten{レ}信\azconvKunten{二}於躬\azconvKunten{一}、而信\azconvKunten{二}於心\azconvKunten{一}。是以難。

\end{azconvIndented}
\begin{azconvIndented}{2\zw}{-1\zw}

〔譯〕\jruby[g]{信}{しん}を人に取るは難し。人は口を信ぜずして\jruby[g]{躬}{み}を信ず。躬を信ぜずして心を信ず。是を以て難し。

\end{azconvIndented}
\begin{azconvIndented}{3\zw}{-1\zw}

評南洲\jruby[g]{守庭吏}{しゆていり}と爲る。島津\jruby[g]{齊彬}{なりあきら}公其の\jruby[g]{眼光}{がんくわう}\jruby[g]{烱々}{けい{〳〵}}として人を\jruby[g]{射}{い}るを見て\jruby[g]{凡}{ぼん}人に非ずと\jruby[g]{以爲}{おも}ひ、\jruby[g]{拔擢}{ばつてき}して之を用ふ。公\jruby[g]{嘗}{かつ}て書を\jruby[g]{作}{つく}り、南洲に命じて之を\jruby[g]{水戸}{みと}の\jruby[g]{烈}{れつ}公に致さしめ、初めより\jruby[g]{封緘}{ふうかん}を加へず。烈公の\jruby[g]{答書}{たふしよ}も亦然り。

\end{azconvIndented}

\begin{azconvIndented}{1\zw}{-1\zw}

二九　臨時之信、累\azconvKunten{二}功於平日\azconvKunten{一}。平日之信、收\azconvKunten{二}効於臨時\azconvKunten{一}。

\end{azconvIndented}
\begin{azconvIndented}{2\zw}{-1\zw}

〔譯〕\jruby[g]{臨時}{りんじ}の\jruby[g]{信}{しん}は、\jruby[g]{功}{こう}を平日に\jruby[g]{累}{かさ}ぬればなり。平日の信は、\jruby[g]{効}{こう}を臨時に\jruby[g]{收}{をさ}むべし。

\end{azconvIndented}
\begin{azconvIndented}{3\zw}{-1\zw}

評南洲官軍の\jruby[g]{先鋒}{せんぱう}となり、品川に\jruby[g]{抵}{いた}る、\jruby[g]{勝安房}{かつあは}、大久保一翁、山岡鐵太郎之を見て、慶喜\jruby[g]{罪}{つみ}を\jruby[g]{俟}{ま}つの\jruby[g]{状}{じやう}を\jruby[g]{具陳}{ぐちん}し、\jruby[g]{討伐}{たうばつ}を\jruby[g]{弛}{ゆる}べんことを請ふ。安房素より南洲を知れり、之を説くこと甚だ力む。乃ち令を諸軍に傳へて、攻撃を\jruby[g]{止}{とゞ}む。

\end{azconvIndented}

\begin{azconvIndented}{1\zw}{-1\zw}

三〇　信孚\azconvKunten{二}於上下\azconvKunten{一}、天下無\azconvKunten{二}甚難\azconvKunten{レ}處事\azconvKunten{一}。

\end{azconvIndented}
\begin{azconvIndented}{2\zw}{-1\zw}

譯信上下に\jruby[g]{孚}{ふ}す、天下甚だ\jruby[g]{處}{しよ}し難き事無し。

\end{azconvIndented}

\begin{azconvIndented}{1\zw}{-1\zw}

三一　意之誠否、須\azconvKunten{下}於\azconvKunten{二}夢寐中事\azconvKunten{一}驗\azconvKunten{\azconvJR}之。

\end{azconvIndented}
\begin{azconvIndented}{2\zw}{-1\zw}

〔譯〕\jruby[g]{意}{い}の\jruby[g]{誠否}{せいひ}は、須らく\jruby[g]{夢寐}{むび}\jruby[g]{中}{ちゆう}の事に於て之を\jruby[g]{驗}{けん}すべし。

\end{azconvIndented}
\begin{azconvIndented}{3\zw}{-1\zw}

評南洲\jruby[g]{弱冠}{じやくくわん}の時、\jruby[g]{藤田東湖}{ふじたとうこ}に\jruby[g]{謁}{えつ}す、東湖は\jruby[g]{重瞳子}{ちやうどうし}、\jruby[g]{躯幹}{くかん}\jruby[g]{魁傑}{くわいけつ}にして、\jruby[g]{黄麻}{わうま}の\jruby[g]{外套}{ぐわいとう}を\jruby[g]{被}{き}、\jruby[g]{朱室}{しゆざや}の\jruby[g]{長劒}{ちやうけん}を\jruby[g]{佩}{さ}して南洲を\jruby[g]{邀}{むか}ふ。南洲一見して\jruby[g]{瞿然}{くぜん}たり。乃ち室内に入る、一大白を\jruby[g]{屬}{ぞく}して\jruby[g]{酒}{さけ}を\jruby[g]{侑}{すゝ}めらる。南洲は\jruby[g]{素}{も}と\jruby[g]{飮}{いん}を\jruby[g]{解}{かい}せず、\jruby[g]{強}{し}ひて之を\jruby[g]{盡}{つく}す、\jruby[g]{忽}{たちま}ち\jruby[g]{酩酊}{めいてい}して\jruby[g]{嘔吐}{おうど}\jruby[g]{席}{せき}を\jruby[g]{汚}{けが}す。東湖は南洲の\jruby[g]{朴率}{ぼくそつ}にして\jruby[g]{飾}{かざ}るところなきを見て\jruby[g]{酷}{はなは}だ之を\jruby[g]{愛}{あい}す。嘗て曰ふ、他日我が志を\jruby[g]{繼}{つ}ぐ者は獨此の少年子のみと。南洲も亦曰ふ、天下\jruby[g]{眞}{しん}に\jruby[g]{畏}{おそ}る可き者なし、\jruby[g]{唯}{たゞ}畏る可き者は東湖一人のみと。二子の言、\jruby[g]{夢寐}{むび}相\jruby[g]{感}{かん}ずる者か。

\end{azconvIndented}

\begin{azconvIndented}{1\zw}{-1\zw}

三二　不\azconvKunten{レ}起\azconvKunten{二}妄念\azconvKunten{一}是敬。妄念不\azconvKunten{レ}起是誠。

\end{azconvIndented}
\begin{azconvIndented}{2\zw}{-1\zw}

〔譯〕\jruby[g]{妄念}{ばうねん}を起さゞるは是れ\jruby[g]{敬}{けい}なり。妄念起らざるは是れ\jruby[g]{誠}{まこと}なり。

\end{azconvIndented}

\begin{azconvIndented}{1\zw}{-1\zw}

三三　因\azconvKunten{二}民義\azconvKunten{一}以激\azconvKunten{レ}之、因\azconvKunten{二}民欲\azconvKunten{一}以趨\azconvKunten{レ}之、則民忘\azconvKunten{二}其生\azconvKunten{一}而致\azconvKunten{二}其死\azconvKunten{一}。是可\azconvKunten{二}以一戰\azconvKunten{一}。

\end{azconvIndented}
\begin{azconvIndented}{2\zw}{-1\zw}

譯民の\jruby[g]{義}{ぎ}に因つて以て之を\jruby[g]{激}{げき}し、民の\jruby[g]{欲}{よく}に因つて以て之を\jruby[g]{趨}{はし}らさば、則ち民其の生を\jruby[g]{忘}{わす}れて其の死を\jruby[g]{致}{いた}さん。是れ以て一\jruby[g]{戰}{せん}す可し。

\end{azconvIndented}
\begin{azconvIndented}{3\zw}{-1\zw}

評兵數は\jruby[g]{孰}{いづ}れか\jruby[g]{衆}{おほ}き、\jruby[g]{器械}{きかい}は孰れか\jruby[g]{精}{せい}なる、\jruby[g]{糧食}{りやうしよく}は孰れか\jruby[g]{積}{つ}める、この數者を以て之を\jruby[g]{較}{くら}べば、\jruby[g]{薩長}{さつちやう}の兵は固より幕府に及ばざるなり。然り而して\jruby[g]{伏見}{ふしみ}の一戰、東兵\jruby[g]{披靡}{ひび}するものは何ぞや。南洲及び木戸公等の\jruby[g]{{※}}{さく}{\small ［＃※は「竹かんむり／束」、41-8］}、民の\jruby[g]{欲}{よく}に因つて之を\jruby[g]{趨}{はし}らしたればなり。是を以て\jruby[g]{破竹}{はちく}の\jruby[g]{勢}{いきほひ}ありたり。

\end{azconvIndented}

\begin{azconvIndented}{1\zw}{-1\zw}

三四　漸必成\azconvKunten{レ}事、惠必懷\azconvKunten{レ}人。如\azconvKunten{二}歴代姦雄\azconvKunten{一}、有\azconvKunten{下}竊\azconvKunten{二}其祕\azconvKunten{一}者\azconvKunten{上}、一時亦能遂\azconvKunten{レ}志。可\azconvKunten{レ}畏之至。

\end{azconvIndented}
\begin{azconvIndented}{2\zw}{-1\zw}

〔譯〕\jruby[g]{漸}{ぜん}は必ず事を\jruby[g]{成}{な}し、\jruby[g]{惠}{けい}は必ず人を\jruby[g]{懷}{な}づく。\jruby[g]{歴代}{れきだい}\jruby[g]{姦雄}{かんゆう}の如き、其\jruby[g]{祕}{ひ}を\jruby[g]{竊}{ぬす}む者有り、一時亦能く志を\jruby[g]{遂}{と}ぐ。畏る可きの至りなり。

\end{azconvIndented}

\begin{azconvIndented}{1\zw}{-1\zw}

三五　匿情似\azconvKunten{二}愼密\azconvKunten{一}。柔媚似\azconvKunten{二}恭順\azconvKunten{一}。剛愎似\azconvKunten{二}自信\azconvKunten{一}。故君子惡\azconvKunten{二}似而非者\azconvKunten{一}。

\end{azconvIndented}
\begin{azconvIndented}{2\zw}{-1\zw}

〔譯〕\jruby[g]{匿情}{とくじやう}は\jruby[g]{愼密}{しんみつ}に\jruby[g]{似}{に}る。\jruby[g]{柔媚}{じうび}は\jruby[g]{恭順}{きようじゆん}に似る。\jruby[g]{剛愎}{がうふく}は\jruby[g]{自信}{じしん}に似る。故に君子は\jruby[g]{似}{に}て\jruby[g]{非}{ひ}なる者を\jruby[g]{惡}{にく}む。

\end{azconvIndented}

\begin{azconvIndented}{1\zw}{-1\zw}

三六　事\azconvKunten{レ}君不\azconvKunten{レ}忠非\azconvKunten{レ}孝也、戰陳無\azconvKunten{レ}勇非\azconvKunten{レ}孝也。曾子孝子、其言如\azconvKunten{レ}此。彼謂\azconvKunten{三}忠孝不\azconvKunten{二}兩全\azconvKunten{一}者、世俗之見也。

\end{azconvIndented}
\begin{azconvIndented}{2\zw}{-1\zw}

譯君に\jruby[g]{事}{つか}へて忠ならざるは孝に非ざるなり、\jruby[g]{戰陳}{せんじん}に\jruby[g]{勇}{ゆう}無きは孝に非ざるなりと。\jruby[g]{曾子}{そうし}は孝子なり、其の言\jruby[g]{此}{かく}の如し。彼の忠孝\jruby[g]{兩全}{りやうぜん}せずと謂ふは、\jruby[g]{世俗}{せぞく}の見なり。

\end{azconvIndented}
\begin{azconvIndented}{3\zw}{-1\zw}

評十年の\jruby[g]{難}{なん}、賊の\jruby[g]{精鋭}{せいえい}熊本城下に\jruby[g]{聚}{あつま}る。而て\jruby[g]{援軍}{えんぐん}未だ達せず。谷中將死を以て之を守り、少しも動かず。\jruby[g]{賊勢}{ぞくせい}遂に屈し、其兵を東する能はず。\jruby[g]{昔者}{むかし}加藤\jruby[g]{嘉明}{よしあき}言へるあり。曰ふ、\jruby[g]{將}{しやう}を\jruby[g]{斬}{き}り\jruby[g]{旗}{はた}を\jruby[g]{搴}{と}るは、氣盛なる者之を能くす、而かも\jruby[g]{眞勇}{しんゆう}に非ざるなり。\jruby[g]{孤城}{こじやう}を\jruby[g]{援}{えん}なきに守り、\jruby[g]{孱}{せん}主を衆\jruby[g]{睽}{そむ}くに\jruby[g]{保}{たも}つ、\jruby[g]{律義者}{りちぎもの}に非ざれば能はず、故に眞勇は必ず\jruby[g]{律義者}{りちぎもの}に出づと。\jruby[g]{尾藤孝肇}{びとうかうてう}曰ふ、\jruby[g]{律義}{りちぎ}とは\jruby[g]{蓋}{けだ}し\jruby[g]{直}{ちよく}にして信あるを謂ふと。余謂ふ、孤城を\jruby[g]{援}{えん}なきに守るは、谷中將の如くば可なりと。嗚呼中將は忠且つ勇なり、而して孝其の\jruby[g]{中}{うち}に在り。

\end{azconvIndented}

\begin{azconvIndented}{1\zw}{-1\zw}

三七　不\azconvKunten{レ}可\azconvKunten{レ}誣者人情、不\azconvKunten{レ}可\azconvKunten{レ}欺者天理、人皆知\azconvKunten{レ}之。蓋知而未\azconvKunten{レ}知。

\end{azconvIndented}
\begin{azconvIndented}{2\zw}{-1\zw}

〔譯〕\jruby[g]{誣}{し}ふ可らざる者は人情なり、\jruby[g]{欺}{あざむ}く可らざる者は天理なり、人皆之を知る。\jruby[g]{蓋}{けだ}し知つて而して未だ知らず。

\end{azconvIndented}
\begin{azconvIndented}{3\zw}{-1\zw}

〔評〕\jruby[g]{榎本武揚}{えのもとぶやう}等五\jruby[g]{稜郭}{りようかく}の兵已に敗る。\jruby[g]{海律全書}{かいりつぜんしよ}二卷を以て我が海軍に\jruby[g]{贈}{おく}つて云ふ、是れ嘗て\jruby[g]{荷蘭}{おらんだ}に學んで\jruby[g]{獲}{え}たる所なり、身と倶に\jruby[g]{滅}{ほろ}ぶることを惜しむと。武揚の誣ふ可らざるの情\jruby[g]{天聽}{てんちやう}に\jruby[g]{達}{たつ}し、其の死を宥し\jruby[g]{寵用}{ちようよう}せらる、天理なり。

\end{azconvIndented}

\begin{azconvIndented}{1\zw}{-1\zw}

三八　知是行之主宰、乾道也。行是知之流行、坤道也。合以成\azconvKunten{二}體躯\azconvKunten{一}。則知行、是二而一、一而二。

\end{azconvIndented}
\begin{azconvIndented}{2\zw}{-1\zw}

〔譯〕\jruby[g]{知}{ち}は是れ\jruby[g]{行}{かう}の\jruby[g]{主宰}{しゆさい}なり、\jruby[g]{乾道}{けんだう}なり。行は是れ知の\jruby[g]{流行}{りうかう}なり、\jruby[g]{坤道}{こんだう}なり。合して以て\jruby[g]{體躯}{たいく}を成す。則ち知行は是れ二にして一、一にして二なり。

\end{azconvIndented}

\begin{azconvIndented}{1\zw}{-1\zw}

三九　學貴\azconvKunten{二}自得\azconvKunten{一}。人徒以\azconvKunten{レ}目讀\azconvKunten{二}有字之書\azconvKunten{一}、故局\azconvKunten{二}於字\azconvKunten{一}、不\azconvKunten{レ}得\azconvKunten{二}通透\azconvKunten{一}。當\azconvKunten{三}以\azconvKunten{レ}心讀\azconvKunten{二}無字之書\azconvKunten{一}、乃洞而有\azconvKunten{二}自得\azconvKunten{一}。

\end{azconvIndented}
\begin{azconvIndented}{2\zw}{-1\zw}

〔譯〕\jruby[g]{學}{がく}は\jruby[g]{自得}{じとく}を\jruby[g]{貴}{たふと}ぶ。人\jruby[g]{徒}{いたづら}に目を以て有字の書を讀む、故に字に\jruby[g]{局}{きよく}し、\jruby[g]{通透}{つうとう}することを得ず。\jruby[g]{當}{まさ}に心を以て無字の書を讀むべし、乃ち\jruby[g]{洞}{とう}して自得するところ有らん。

\end{azconvIndented}

\begin{azconvIndented}{1\zw}{-1\zw}

四〇　孟子以\azconvKunten{二}讀書\azconvKunten{一}爲\azconvKunten{二}尚友\azconvKunten{一}。故讀\azconvKunten{二}經籍\azconvKunten{一}、即是聽\azconvKunten{二}嚴師父兄之訓\azconvKunten{一}也。讀\azconvKunten{二}史子\azconvKunten{一}、亦即與\azconvKunten{二}明君賢相英雄豪傑\azconvKunten{一}相周旋也。其可\azconvKunten{レ}不\azconvKunten{下}清\azconvKunten{二}明其心\azconvKunten{一}以對\azconvKunten{中}越之\azconvKunten{上}乎。

\end{azconvIndented}
\begin{azconvIndented}{2\zw}{-1\zw}

譯孟子讀書を以て\jruby[g]{尚友}{しやういう}と爲す。故に\jruby[g]{經籍}{けいせき}を讀む、即ち是れ\jruby[g]{嚴師}{げんし}父兄の訓を聽くなり。\jruby[g]{史子}{しゝ}を讀む、亦即ち明君賢相英雄豪傑と相\jruby[g]{周旋}{しうせん}するなり。其れ其の心を清明にして以て之に\jruby[g]{對越}{たいえつ}せざる可けんや。

\end{azconvIndented}

\begin{azconvIndented}{1\zw}{-1\zw}

四一　爲\azconvKunten{レ}學緊要、在\azconvKunten{二}心一字\azconvKunten{一}。把\azconvKunten{レ}心以治\azconvKunten{レ}心、謂\azconvKunten{二}之聖學\azconvKunten{一}。爲\azconvKunten{レ}政著眼、在\azconvKunten{二}情一字\azconvKunten{一}。循\azconvKunten{レ}情以治\azconvKunten{レ}情、謂\azconvKunten{二}之王道\azconvKunten{一}。王道聖學非\azconvKunten{レ}二。

\end{azconvIndented}
\begin{azconvIndented}{2\zw}{-1\zw}

譯學を爲すの\jruby[g]{緊要}{きんえう}は心の一字に在り。心を\jruby[g]{把}{と}つて以て心を治む、之を聖學と謂ふ。政を爲すの\jruby[g]{着眼}{ちやくがん}は情の一字に在り。情に\jruby[g]{循}{したが}うて以て情を治む、之を王道と謂ふ。王道と聖學と二に非ず。

\end{azconvIndented}
\begin{azconvIndented}{3\zw}{-1\zw}

評兵を\jruby[g]{治}{ち}して\jruby[g]{對抗}{たいかう}し、互に\jruby[g]{勝敗}{しようはい}あり。兵士或は\jruby[g]{負傷}{ふしやう}者の\jruby[g]{状}{じやう}を爲す、\jruby[g]{醫}{い}故に之を\jruby[g]{診察}{しんさつ}す。兵士初め負傷者とならんことを惡む。一日、\jruby[g]{聖上}{せいじやう}\jruby[g]{親臨}{しんりん}して負傷者を\jruby[g]{撫}{ぶ}し、\jruby[g]{恩言}{おんげん}を\jruby[g]{賜}{たま}ふ、此より兵士負傷者とならんことを願ふ。是に由つて之を觀れば、兵を\jruby[g]{馭}{ぎよ}するも亦情に外ならざるなり。

\end{azconvIndented}

\begin{azconvIndented}{1\zw}{-1\zw}

四二　發\azconvKunten{レ}憤忘\azconvKunten{レ}食、志氣如\azconvKunten{レ}是。樂以忘\azconvKunten{レ}憂、心體如\azconvKunten{レ}是。不\azconvKunten{レ}知\azconvKunten{二}老之將\azconvKunten{\azconvIR}至、知\azconvKunten{レ}命樂\azconvKunten{レ}天如\azconvKunten{レ}是。聖人與\azconvKunten{レ}人不\azconvKunten{レ}同、又與\azconvKunten{レ}人不\azconvKunten{レ}異。

\end{azconvIndented}
\begin{azconvIndented}{2\zw}{-1\zw}

〔譯〕\jruby[g]{憤}{いきどほり}を發して食を\jruby[g]{忘}{わす}る、\jruby[g]{志氣}{しき}\jruby[g]{是}{かく}の如し。\jruby[g]{樂}{たのし}んで以て\jruby[g]{憂}{うれひ}を忘る、\jruby[g]{心體}{しんたい}是の如し。\jruby[g]{老}{らう}の將に至らんとするを知らず、\jruby[g]{命}{めい}を知り天を樂しむもの\jruby[g]{是}{かく}の如し。聖人は人と同じからず、又人と\jruby[g]{異}{こと}ならず。

\end{azconvIndented}

\begin{azconvIndented}{1\zw}{-1\zw}

四三　講\azconvKunten{二}説聖賢\azconvKunten{一}、而不\azconvKunten{レ}能\azconvKunten{レ}躬\azconvKunten{レ}之、謂\azconvKunten{二}之口頭聖賢\azconvKunten{一}、吾聞\azconvKunten{レ}之一惕然。論\azconvKunten{二}辯道學\azconvKunten{一}、而不\azconvKunten{レ}能\azconvKunten{レ}體\azconvKunten{レ}之、謂\azconvKunten{二}之紙上道學\azconvKunten{一}、吾聞\azconvKunten{レ}之再惕然。

\end{azconvIndented}
\begin{azconvIndented}{2\zw}{-1\zw}

譯聖賢を\jruby[g]{講説}{かうせつ}して之を\jruby[g]{躬}{み}にする能はず、之を\jruby[g]{口頭}{こうとう}聖賢と謂ふ、吾れ之を聞いて一たび\jruby[g]{惕然}{てきぜん}たり。道學を\jruby[g]{論辯}{ろんべん}して之を\jruby[g]{體}{たい}する能はず、之を紙上道學と謂ふ、吾れ之を聞いて再び\jruby[g]{惕然}{てきぜん}たり。

\end{azconvIndented}

\begin{azconvIndented}{1\zw}{-1\zw}

四四　學、稽\azconvKunten{二}之古訓\azconvKunten{一}、問、質\azconvKunten{二}之師友\azconvKunten{一}、人皆知\azconvKunten{レ}之。學必學\azconvKunten{二}之躬\azconvKunten{一}、問必問\azconvKunten{二}諸心\azconvKunten{一}、其有\azconvKunten{二}幾人\azconvKunten{一}耶。

\end{azconvIndented}
\begin{azconvIndented}{2\zw}{-1\zw}

〔譯〕\jruby[g]{學}{がく}之を\jruby[g]{古訓}{こくん}に\jruby[g]{稽}{かんが}へ、\jruby[g]{問}{もん}之を師友に\jruby[g]{質}{たゞ}すは、人皆之を知る。學必ず之を躬に學び、問必ず諸を心に問ふは、其れ幾人有らんか。

\end{azconvIndented}

\begin{azconvIndented}{1\zw}{-1\zw}

四五　以\azconvKunten{レ}天而得者固。以\azconvKunten{レ}人而得者脆。

\end{azconvIndented}
\begin{azconvIndented}{2\zw}{-1\zw}

譯天を以て得たるものは\jruby[g]{固}{かた}し。人を以て得たるものは\jruby[g]{脆}{もろ}し。

\end{azconvIndented}

\begin{azconvIndented}{1\zw}{-1\zw}

四六　君子自慊、小人自欺。君子自彊、小人自棄。上達下達、落\azconvKunten{二}在一自字\azconvKunten{一}。

\end{azconvIndented}
\begin{azconvIndented}{2\zw}{-1\zw}

譯君子は自ら\jruby[g]{慊}{こゝろよ}くし、小人は自ら\jruby[g]{欺}{あざむ}く。君子は自ら\jruby[g]{彊}{つと}め、小人は自ら\jruby[g]{棄}{す}つ。上\jruby[g]{達}{たつ}と下\jruby[g]{達}{たつ}とは、一の\jruby[g]{自}{じ}の字に\jruby[g]{落在}{らくざい}す。

\end{azconvIndented}

\begin{azconvIndented}{1\zw}{-1\zw}

四七　人皆知\azconvKunten{レ}問\azconvKunten{二}身之安否\azconvKunten{一}、而不\azconvKunten{レ}知\azconvKunten{レ}問\azconvKunten{二}心之安否\azconvKunten{一}。宜\azconvKunten{下}自問\azconvKunten{中}能不\azconvKunten{レ}欺\azconvKunten{二}闇室\azconvKunten{一}否、能不\azconvKunten{レ}愧\azconvKunten{二}衾影\azconvKunten{一}否、能得\azconvKunten{二}安穩快樂\azconvKunten{一}否\azconvKunten{上}。時時如\azconvKunten{レ}是、心便不\azconvKunten{レ}放。

\end{azconvIndented}
\begin{azconvIndented}{2\zw}{-1\zw}

譯人は皆身の\jruby[g]{安否}{あんぴ}を\jruby[g]{問}{と}ふことを知つて、而かも心の安否を問ふことを知らず。宜しく自ら能く\jruby[g]{闇室}{あんしつ}を\jruby[g]{欺}{あざむ}かざるや\jruby[g]{否}{いな}や、能く\jruby[g]{衾影}{きんえい}に\jruby[g]{愧}{は}ぢざるや否や、能く\jruby[g]{安穩}{あんおん}\jruby[g]{快樂}{くわいらく}を得るや否やと問ふべし。時時\jruby[g]{是}{かく}の如くば心\jruby[g]{便}{すなは}ち\jruby[g]{放}{はな}たず。

\end{azconvIndented}
\begin{azconvIndented}{3\zw}{-1\zw}

評某士南洲に\jruby[g]{面}{めん}して\jruby[g]{仕官}{しくわん}を\jruby[g]{求}{もと}む。南洲曰ふ、汝\jruby[g]{俸給}{ほうきふ}\jruby[g]{幾許}{いくばく}を求むるやと。某曰ふ、三十圓ばかりと。南洲乃ち三十圓を與へて曰ふ、汝に\jruby[g]{一月}{ひとつき}の\jruby[g]{俸}{ほう}金を與へん、汝は宜しく汝の心に\jruby[g]{向}{むか}うて我が\jruby[g]{才力}{さいりき}如何を問ふべしと。其人\jruby[g]{復}{ま}た來らず。

\end{azconvIndented}

\begin{azconvIndented}{1\zw}{-1\zw}

四八　無\azconvKunten{レ}爲而有\azconvKunten{レ}爲之謂\azconvKunten{レ}誠。有\azconvKunten{レ}爲而無\azconvKunten{レ}爲之謂\azconvKunten{レ}敬。

\end{azconvIndented}
\begin{azconvIndented}{2\zw}{-1\zw}

譯爲す無くして爲す有る之を\jruby[g]{誠}{まこと}と謂ふ。爲す有つて爲す無し之を\jruby[g]{敬}{けい}と謂ふ。

\end{azconvIndented}

\begin{azconvIndented}{1\zw}{-1\zw}

四九　寛懷不\azconvKunten{レ}忤\azconvKunten{二}俗情\azconvKunten{一}、和也。立脚不\azconvKunten{レ}墜\azconvKunten{二}俗情\azconvKunten{一}、介也。

\end{azconvIndented}
\begin{azconvIndented}{2\zw}{-1\zw}

〔譯〕\jruby[g]{寛懷}{かんくわい}\jruby[g]{俗情}{ぞくじやう}に\jruby[g]{忤}{さか}はざるは、\jruby[g]{和}{わ}なり。\jruby[g]{立脚}{りつきやく}俗情に\jruby[g]{墜}{お}ちざるは、\jruby[g]{介}{かい}なり。

\end{azconvIndented}

\begin{azconvIndented}{1\zw}{-1\zw}

五〇　惻隱之心偏、民或有\azconvKunten{二}溺\azconvKunten{レ}愛殞\azconvKunten{レ}身者\azconvKunten{一}。羞惡之心偏、民或有\azconvKunten{下}自\azconvKunten{二}經溝涜\azconvKunten{一}者\azconvKunten{上}。辭讓之心偏、民或有\azconvKunten{二}奔亡風狂者\azconvKunten{一}。是非之心偏、民或有\azconvKunten{二}兄弟鬩\azconvKunten{レ}牆父子相訟者\azconvKunten{一}。凡情之偏、雖\azconvKunten{二}四端\azconvKunten{一}遂陷\azconvKunten{二}不善\azconvKunten{一}。故學以致\azconvKunten{二}中和\azconvKunten{一}、歸\azconvKunten{三}於無\azconvKunten{二}過不及\azconvKunten{一}、謂\azconvKunten{二}之復性之學\azconvKunten{一}。

\end{azconvIndented}
\begin{azconvIndented}{2\zw}{-1\zw}

〔譯〕\jruby[g]{惻隱}{そくいん}の心\jruby[g]{偏}{へん}すれば、民或は\jruby[g]{愛}{あい}に\jruby[g]{溺}{おぼ}れ身を\jruby[g]{殞}{おと}す者有り。\jruby[g]{羞惡}{しうを}の心偏すれば、民或は\jruby[g]{溝涜}{かうとく}に\jruby[g]{自經}{じけい}する者有り。\jruby[g]{辭讓}{じじやう}の心偏すれば、民或は\jruby[g]{奔亡}{ほんばう}\jruby[g]{風狂}{ふうきやう}する者有り。是非の心偏すれば、民或は兄弟\jruby[g]{牆}{かき}に\jruby[g]{鬩}{せめ}ぎ父子相\jruby[g]{訟}{うつた}ふ者有り。凡そ情の偏するや、四\jruby[g]{端}{たん}と雖遂に\jruby[g]{不善}{ふぜん}に\jruby[g]{陷}{おちい}る。故に學んで以て中和を\jruby[g]{致}{いた}し、\jruby[g]{過不及}{かふきふ}無きに\jruby[g]{歸}{き}す、之を\jruby[g]{復性}{ふくせい}の學と謂ふ。

\end{azconvIndented}
\begin{azconvIndented}{3\zw}{-1\zw}

評江藤\jruby[g]{新平}{しんぺい}、前原\jruby[g]{一誠}{いつせい}等の如きは、皆\jruby[g]{維新}{いしん}の功臣として、勤王二なく、官は\jruby[g]{參議}{さんぎ}に至り、位は人臣の\jruby[g]{榮}{えい}を\jruby[g]{極}{きは}む。然り而して前後皆亂を爲し誅に伏す、惜しいかな。豈四\jruby[g]{端}{たん}の\jruby[g]{偏}{へん}ありしものか。

\end{azconvIndented}

\begin{azconvIndented}{1\zw}{-1\zw}

五一　此學吾人一生負擔、當\azconvKunten{二}斃而後已\azconvKunten{一}。道固無\azconvKunten{レ}窮、堯舜之上善無\azconvKunten{レ}盡。孔子自\azconvKunten{レ}志\azconvKunten{レ}學、至\azconvKunten{二}七十\azconvKunten{一}、毎\azconvKunten{二}十年\azconvKunten{一}、自覺\azconvKunten{二}其有\azconvKunten{\azconvIR}所\azconvKunten{レ}進、孜孜自彊、不\azconvKunten{レ}知\azconvKunten{二}老之將\azconvKunten{\azconvIR}至。假使\azconvKunten{二}其踰\azconvKunten{レ}耄至\azconvKunten{\azconvIR}期、則其神明不\azconvKunten{レ}測、想當\azconvKunten{レ}爲\azconvKunten{二}何如\azconvKunten{一}哉。凡學\azconvKunten{二}孔子\azconvKunten{一}者、宜\azconvKunten{下}以\azconvKunten{二}孔子之志\azconvKunten{一}爲\azconvKunten{\azconvJR}志。

\end{azconvIndented}
\begin{azconvIndented}{2\zw}{-1\zw}

譯此の學は吾人一生の\jruby[g]{負擔}{ふたん}、\jruby[g]{當}{まさ}に\jruby[g]{斃}{たふ}れて後に\jruby[g]{已}{や}むべし。道固より窮り無し。堯舜の上、善盡くること無し。孔子學に志してより七十に至るまで、十年毎に自ら其の\jruby[g]{進}{すゝ}む所有るを\jruby[g]{覺}{さと}り、\jruby[g]{孜孜}{しゝ}として自ら\jruby[g]{彊}{つと}めて、\jruby[g]{老}{らう}の將に至らんとするを知らず。\jruby[g]{假}{も}し其をして\jruby[g]{耄}{ばう}を\jruby[g]{踰}{こ}え\jruby[g]{期}{き}に至らしめば、則ち其の神明\jruby[g]{測}{はか}られざること、\jruby[g]{想}{おも}ふに當に何如たるべきぞや。凡そ孔子を學ぶ者は、宜しく孔子の志を以て志と爲すべし。

\end{azconvIndented}

\begin{azconvIndented}{1\zw}{-1\zw}

五二　自彊不\azconvKunten{レ}息、天道也、君子所\azconvKunten{レ}以也。如\azconvKunten{下}虞舜孳孳爲\azconvKunten{レ}善、大禹思\azconvKunten{二}日孜孜\azconvKunten{一}、成湯苟日新、文王不\azconvKunten{二}遑暇\azconvKunten{一}、周公坐以待\azconvKunten{レ}旦、孔子發\azconvKunten{レ}憤忘\azconvKunten{\azconvJR}食、皆是也。彼徒事\azconvKunten{二}靜養瞑坐\azconvKunten{一}而已、則與\azconvKunten{二}此學脈\azconvKunten{一}背馳。

\end{azconvIndented}
\begin{azconvIndented}{2\zw}{-1\zw}

譯自ら\jruby[g]{彊}{つと}めて\jruby[g]{息}{や}まざるは天道なり、君子の\jruby[g]{以}{もち}ゐる所なり。\jruby[g]{虞舜}{ぐしゆん}の\jruby[g]{孳孳}{じじ}として善を爲し、大\jruby[g]{禹}{う}の日に孜孜せんことを思ひ、\jruby[g]{成湯}{せいたう}の\jruby[g]{苟}{まこと}に日に新にせる、文王の\jruby[g]{遑}{いとま}あき\jruby[g]{暇}{いとま}あらざる、\jruby[g]{周}{しう}公の\jruby[g]{坐}{ざ}して以て\jruby[g]{旦}{たん}を\jruby[g]{待}{ま}つ、孔子の\jruby[g]{憤}{いきどほ}りを發して食を忘るゝ如きは、皆是なり。彼の\jruby[g]{徒}{いたづら}に\jruby[g]{靜養}{せいやう}\jruby[g]{瞑坐}{めいざ}を事とすのみならば、則ち此の\jruby[g]{學脈}{がくみやく}と\jruby[g]{背馳}{はいち}す。

\end{azconvIndented}

\begin{azconvIndented}{1\zw}{-1\zw}

五三　自彊不\azconvKunten{レ}息時候、心地光光明明、有\azconvKunten{二}何妄念游思\azconvKunten{一}、有\azconvKunten{二}何嬰累罣想\azconvKunten{一}。

\end{azconvIndented}
\begin{azconvIndented}{2\zw}{-1\zw}

譯自ら\jruby[g]{彊}{つと}めて\jruby[g]{息}{や}まざる\jruby[g]{時候}{じこう}は、\jruby[g]{心地}{しんち}\jruby[g]{光光明明}{くわう{〳〵}めい{〳〵}}にして、何の\jruby[g]{妄念}{ばうねん}\jruby[g]{游思}{ゆうし}有らん、何の\jruby[g]{嬰累}{えいるゐ}\jruby[g]{罣想}{けさう}有らん。

\end{azconvIndented}
\begin{azconvIndented}{3\zw}{-1\zw}

評三條公の筑前に在る、或る人其の\jruby[g]{旅況}{りよきやう}の\jruby[g]{無聊}{むれう}を\jruby[g]{察}{さつ}して美女を進む、公之を\jruby[g]{卻}{しりぞ}く。某氏\jruby[g]{宴}{えん}を\jruby[g]{開}{ひら}いて女\jruby[g]{樂}{がく}を\jruby[g]{設}{まう}く、公\jruby[g]{怫}{ふつ}然として去れり。

\end{azconvIndented}

\begin{azconvIndented}{1\zw}{-1\zw}

五四　提\azconvKunten{二}一燈\azconvKunten{一}、行\azconvKunten{二}暗夜\azconvKunten{一}。勿\azconvKunten{レ}憂\azconvKunten{二}暗夜\azconvKunten{一}、只頼\azconvKunten{二}一燈\azconvKunten{一}。

\end{azconvIndented}
\begin{azconvIndented}{2\zw}{-1\zw}

譯一\jruby[g]{燈}{とう}を\jruby[g]{提}{ひつさ}げて、\jruby[g]{暗夜}{あんや}を行く。暗夜を\jruby[g]{憂}{うれ}ふる勿れ、只だ一\jruby[g]{燈}{とう}を\jruby[g]{頼}{たの}め。

\end{azconvIndented}
\begin{azconvIndented}{3\zw}{-1\zw}

〔評〕\jruby[g]{伏水}{ふしみ}戰を開き、\jruby[g]{砲聲}{はうせい}\jruby[g]{大内}{おほうち}に聞え、愈\jruby[g]{激}{はげ}しく愈\jruby[g]{近}{ちか}づく。岩倉公南洲に問うて曰ふ、\jruby[g]{勝敗}{しようはい}何如と。南洲答へて曰ふ、西郷隆盛在り、憂ふる勿れと。

\end{azconvIndented}

\begin{azconvIndented}{1\zw}{-1\zw}

五五　倫理物理、同一理也。我學\azconvKunten{二}倫理之學\azconvKunten{一}、宜\azconvKunten{三}近取\azconvKunten{二}諸身\azconvKunten{一}、即是物理。

\end{azconvIndented}
\begin{azconvIndented}{2\zw}{-1\zw}

〔譯〕\jruby[g]{倫理}{りんり}と物理とは同一理なり。我れ倫理の學を學ぶ、宜しく近く\jruby[g]{諸}{これ}を身に取るべし、即ち是れ物理なり。

\end{azconvIndented}

\begin{azconvIndented}{1\zw}{-1\zw}

五六　濁水亦水也。一澄則爲\azconvKunten{二}清水\azconvKunten{一}。客氣亦氣也。一轉則爲\azconvKunten{二}正氣\azconvKunten{一}。逐\azconvKunten{レ}客工夫、只是克\azconvKunten{レ}己、只是復\azconvKunten{レ}禮。

\end{azconvIndented}
\begin{azconvIndented}{2\zw}{-1\zw}

〔譯〕\jruby[g]{濁水}{だくすゐ}も亦水なり、一\jruby[g]{澄}{ちよう}すれば則ち\jruby[g]{清水}{せいすゐ}となる。\jruby[g]{客氣}{きやくき}も亦氣なり、一\jruby[g]{轉}{てん}すれば則ち\jruby[g]{正氣}{せいき}となる。\jruby[g]{客}{きやく}を\jruby[g]{逐}{お}ふの工夫は、只是れ己に克つなり、只是れ禮に\jruby[g]{復}{かへ}るなり。

\end{azconvIndented}
\begin{azconvIndented}{3\zw}{-1\zw}

評南洲\jruby[g]{壯時}{さうじ}\jruby[g]{角觝}{かくてい}を好み、\jruby[g]{毎}{つね}に壯士と角す。人之を\jruby[g]{苦}{くる}しむ。其\jruby[g]{守庭吏}{しゆていり}と爲るや、\jruby[g]{庭}{てい}中に\jruby[g]{土豚}{どとん}を\jruby[g]{設}{まう}けて、\jruby[g]{掃除}{さうぢよ}を\jruby[g]{事}{こと}とせず。既にして\jruby[g]{慨然}{がいぜん}として天下を以て自ら\jruby[g]{任}{にん}じ、\jruby[g]{節}{せつ}を\jruby[g]{屈}{くつ}して書を讀み、遂に\jruby[g]{復古}{ふくこ}の大\jruby[g]{業}{げふ}を成せり。

\end{azconvIndented}

\begin{azconvIndented}{1\zw}{-1\zw}

五七　理本無\azconvKunten{レ}形。無\azconvKunten{レ}形則無\azconvKunten{レ}名矣。形而後有\azconvKunten{レ}名。既有\azconvKunten{レ}名、則理謂\azconvKunten{二}之氣\azconvKunten{一}無\azconvKunten{二}不可\azconvKunten{一}。故專指\azconvKunten{二}本體\azconvKunten{一}、則形後亦謂\azconvKunten{二}之理\azconvKunten{一}。專指\azconvKunten{二}運用\azconvKunten{一}、則形前亦謂\azconvKunten{二}之氣\azconvKunten{一}、竝無\azconvKunten{二}不可\azconvKunten{一}。如\azconvKunten{二}浩然之氣\azconvKunten{一}、專指\azconvKunten{二}運用\azconvKunten{一}、其實太極之呼吸、只是一誠。謂\azconvKunten{二}之氣原\azconvKunten{一}、即是理。

\end{azconvIndented}
\begin{azconvIndented}{2\zw}{-1\zw}

譯理は\jruby[g]{本}{も}と\jruby[g]{形}{かたち}無し。形無ければ則ち名無し。形ありて後に名有り。既に名有れば、則ち理之を氣と謂ふも、不可無し。故に專ら\jruby[g]{本體}{ほんたい}を指せば、則ち\jruby[g]{形後}{けいご}も亦之を理と謂ふ。專ら\jruby[g]{運用}{うんよう}を指せば、則ち形前も亦之を氣と謂ふ、\jruby[g]{竝}{ならび}に不可無し。\jruby[g]{浩然}{かうぜん}の氣の如きは、專ら運用を指すも、其の實\jruby[g]{太極}{たいきよく}の\jruby[g]{呼吸}{こきふ}にして、只是れ一\jruby[g]{誠}{せい}なり。之を氣\jruby[g]{原}{げん}と謂ふ、即ち是れ理なり。

\end{azconvIndented}

\begin{azconvIndented}{1\zw}{-1\zw}

五八　物我一體、即是仁。我執\azconvKunten{二}公情\azconvKunten{一}以行\azconvKunten{二}公事\azconvKunten{一}、天下無\azconvKunten{レ}不\azconvKunten{レ}服。治亂之機、在\azconvKunten{二}於公不公\azconvKunten{一}。周子曰、公\azconvKunten{二}於己\azconvKunten{一}者、公\azconvKunten{二}於人\azconvKunten{一}。伊川又以\azconvKunten{二}公理\azconvKunten{一}、釋\azconvKunten{二}仁字\azconvKunten{一}。餘姚亦更\azconvKunten{二}博愛\azconvKunten{一}爲\azconvKunten{二}公愛\azconvKunten{一}。可\azconvKunten{二}并攷\azconvKunten{一}。

\end{azconvIndented}
\begin{azconvIndented}{2\zw}{-1\zw}

〔譯〕\jruby[g]{物我}{ぶつが}一\jruby[g]{體}{たい}は即ち是れ仁なり。我れ\jruby[g]{公情}{こうじやう}を\jruby[g]{執}{と}つて以て公事を行ふ、天下服せざる無し。\jruby[g]{治亂}{ちらん}の\jruby[g]{機}{き}は公と不公とに在り。\jruby[g]{周}{しう}子曰ふ、\jruby[g]{己}{おのれ}に公なる者は人に公なりと。\jruby[g]{伊川}{いせん}又\jruby[g]{公理}{こうり}を以て仁の字を\jruby[g]{釋}{しやく}す。\jruby[g]{餘姚}{よえう}も亦博愛を\jruby[g]{更}{あらた}めて公愛と爲せり。\jruby[g]{并}{あは}せ\jruby[g]{攷}{かんが}ふ可し。

\end{azconvIndented}
\begin{azconvIndented}{3\zw}{-1\zw}

評余嘗て木戸公の言を記せり。曰ふ、\jruby[g]{會津藩士}{あひづはんし}は、性直にして用ふ可し、\jruby[g]{長人}{ちやうじん}の及ぶ所に非ざるなりと。夫れ\jruby[g]{會}{くわい}は\jruby[g]{長}{ちやう}の\jruby[g]{敵}{てき}なり、\jruby[g]{而}{し}かも其の言\jruby[g]{此}{かく}の如し。以て公の事を\jruby[g]{處}{しよ}すること皆\jruby[g]{公平}{こうへい}なるを知るべし。

\end{azconvIndented}

\begin{azconvIndented}{1\zw}{-1\zw}

五九　尊\azconvKunten{二}徳性\azconvKunten{一}、是以道\azconvKunten{二}問學\azconvKunten{一}、即是尊\azconvKunten{二}徳性\azconvKunten{一}。先立\azconvKunten{二}其大者\azconvKunten{一}、則其知也眞。能迪\azconvKunten{二}其知\azconvKunten{一}、則其功也實。畢竟一條路往來耳。

\end{azconvIndented}
\begin{azconvIndented}{2\zw}{-1\zw}

譯徳性を尊ぶ、是を以て\jruby[g]{問學}{ぶんがく}に\jruby[g]{道}{よ}る、即ち是れ徳性を尊ぶなり。先づ其の大なる者を立つれば、則ち其知や\jruby[g]{眞}{しん}なり。能く其の知を\jruby[g]{迪}{ふ}めば、則ち其功や\jruby[g]{實}{じつ}なり。\jruby[g]{畢竟}{ひつきやう}\jruby[g]{一條}{いちでう}\jruby[g]{路}{ろ}の往來のみ。

\end{azconvIndented}

\begin{azconvIndented}{1\zw}{-1\zw}

六〇　周子主\azconvKunten{レ}靜、謂\azconvKunten{三}心守\azconvKunten{二}本體\azconvKunten{一}。※{\small ［＃※は「圖」の「回」に代えて「面から一、二画目をとったもの」、52-8］}説自\azconvKunten{二}註無\azconvKunten{レ}欲故靜\azconvKunten{一}、程伯氏因\azconvKunten{レ}此有\azconvKunten{二}天理人欲之説\azconvKunten{一}。叔子持\azconvKunten{レ}敬工夫亦在\azconvKunten{レ}此。朱陸以下雖\azconvKunten{三}各有\azconvKunten{二}得\azconvKunten{レ}力處\azconvKunten{一}、而畢竟不\azconvKunten{レ}出\azconvKunten{二}此範圍\azconvKunten{一}。不\azconvKunten{レ}意至\azconvKunten{二}明儒\azconvKunten{一}、朱陸分\azconvKunten{レ}黨如\azconvKunten{二}敵讐\azconvKunten{一}。何以然邪。今之學者、宜\azconvKunten{下}以\azconvKunten{二}平心\azconvKunten{一}待\azconvKunten{\azconvJR}之。取\azconvKunten{二}其得\azconvKunten{レ}力處\azconvKunten{一}可也。

\end{azconvIndented}
\begin{azconvIndented}{2\zw}{-1\zw}

〔譯〕\jruby[g]{周子}{しうし}\jruby[g]{靜}{せい}を\jruby[g]{主}{しゆ}とす、\jruby[g]{心}{こゝろ}\jruby[g]{本體}{ほんたい}を守るを謂ふなり。\jruby[g]{{※説}}{づせつ}{\small ［＃※は「圖」の「回」に代えて「面から一、二画目をとったもの」、52-12］}に、「\jruby[g]{欲}{よく}無し故に\jruby[g]{靜}{せい}」と\jruby[g]{自註}{じちゆう}す、\jruby[g]{程伯氏}{ていはくし}\jruby[g]{此}{これ}に因つて天\jruby[g]{理}{り}人\jruby[g]{欲}{よく}の\jruby[g]{説}{せつ}有り。\jruby[g]{叔子}{しゆくし}\jruby[g]{敬}{けい}を\jruby[g]{持}{ぢ}する\jruby[g]{工夫}{くふう}も亦\jruby[g]{此}{こゝ}に在り。\jruby[g]{朱陸}{しゆりく}以下各\jruby[g]{力}{ちから}を得る處有りと雖、\jruby[g]{而}{し}かも\jruby[g]{畢竟}{ひつきやう}此の\jruby[g]{範圍}{はんい}を出でず。\jruby[g]{意}{おも}はざりき\jruby[g]{明儒}{みんじゆ}に至つて、\jruby[g]{朱陸}{しゆりく}\jruby[g]{黨}{たう}を分つこと\jruby[g]{敵讐}{てきしう}の如くあらんとは。何を以て然るや。今の學ぶ者、宜しく平心を以て之を待つべし。其の力を得る處を取らば可なり。

\end{azconvIndented}

\begin{azconvIndented}{1\zw}{-1\zw}

六一　象山、宇宙内事、皆己分内事、此謂\azconvKunten{二}男子擔當之志如\azconvKunten{\azconvIR}此。陳澔引\azconvKunten{レ}此註\azconvKunten{二}射義\azconvKunten{一}、極是。

\end{azconvIndented}
\begin{azconvIndented}{2\zw}{-1\zw}

〔譯〕\jruby[g]{象山}{しようざん}の、\jruby[g]{宇宙}{うちう}\jruby[g]{内}{ない}の事は皆\jruby[g]{己}{おの}れ\jruby[g]{分内}{ぶんない}の事は、\jruby[g]{此}{こ}れ男子\jruby[g]{擔當}{たんたう}の志\jruby[g]{此}{かく}の如きを謂ふなり。\jruby[g]{陳澔}{ちんかう}此を引いて\jruby[g]{射義}{しやぎ}を\jruby[g]{註}{ちゆう}す、\jruby[g]{極}{きは}めて\jruby[g]{是}{ぜ}なり。

\end{azconvIndented}
\begin{azconvIndented}{3\zw}{-1\zw}

評南洲\jruby[g]{嘗}{かつ}て東湖に從うて學ぶ。\jruby[g]{當時}{たうじ}書する所、今猶民間に\jruby[g]{存}{そん}す。曰ふ、「\jruby[g]{一寸}{いつすん}の\jruby[g]{英心}{えいしん}\jruby[g]{萬夫}{ばんぷ}に\jruby[g]{敵}{てき}す」と。\jruby[g]{蓋}{けだ}し\jruby[g]{復古}{ふくこ}の\jruby[g]{業}{げふ}を以て\jruby[g]{擔當}{たんたう}することを爲す。\jruby[g]{維新}{いしん}征東の\jruby[g]{功}{こう}實に此に\jruby[g]{讖}{しん}す。\jruby[g]{末路}{まつろ}\jruby[g]{再}{ふたゝ}び\jruby[g]{讖}{しん}を成せるは、\jruby[g]{悲}{かな}しむべきかな。

\end{azconvIndented}

\begin{azconvIndented}{1\zw}{-1\zw}

六二　講\azconvKunten{二}論語\azconvKunten{一}、是慈父教\azconvKunten{レ}子意思。講\azconvKunten{二}孟子\azconvKunten{一}、是伯兄誨\azconvKunten{レ}季意思。講\azconvKunten{二}大學\azconvKunten{一}、如\azconvKunten{二}網在\azconvKunten{\azconvIR}綱。講\azconvKunten{二}中庸\azconvKunten{一}、如\azconvKunten{二}雲出\azconvKunten{\azconvIR}岫。

\end{azconvIndented}
\begin{azconvIndented}{2\zw}{-1\zw}

〔譯〕\jruby[g]{論語}{ろんご}を\jruby[g]{講}{かう}ず、是れ\jruby[g]{慈父}{じふ}の子を教ふる\jruby[g]{意思}{いし}。\jruby[g]{孟子}{まうし}を講ず、是れ伯兄の\jruby[g]{季}{き}を\jruby[g]{誨}{をし}ふる\jruby[g]{意思}{いし}。\jruby[g]{大學}{だいがく}を講ず、\jruby[g]{網}{あみ}の\jruby[g]{綱}{かう}に在る如し。\jruby[g]{中庸}{ちゆうよう}を講ず、\jruby[g]{雲}{くも}の\jruby[g]{岫}{しう}を出づる如し。

\end{azconvIndented}

\begin{azconvIndented}{1\zw}{-1\zw}

六三　易是性字註脚。詩是情字註脚。書是心字註脚。

\end{azconvIndented}
\begin{azconvIndented}{2\zw}{-1\zw}

〔譯〕\jruby[g]{易}{えき}は是れ\jruby[g]{性}{せい}の字の\jruby[g]{註脚}{ちゆうきやく}なり。\jruby[g]{詩}{し}は是れ情の字の註脚なり。\jruby[g]{書}{しよ}は是れ心の字の註脚なり。

\end{azconvIndented}

\begin{azconvIndented}{1\zw}{-1\zw}

六四　獨得之見似\azconvKunten{レ}私、人驚\azconvKunten{二}其驟至\azconvKunten{一}。平凡之議似\azconvKunten{レ}公、世安\azconvKunten{二}其狃聞\azconvKunten{一}。凡聽\azconvKunten{二}人言\azconvKunten{一}、宜\azconvKunten{二}虚懷而邀\azconvKunten{\azconvIR}之。勿\azconvKunten{レ}苟\azconvKunten{二}安狃聞\azconvKunten{一}可也。

\end{azconvIndented}
\begin{azconvIndented}{2\zw}{-1\zw}

〔譯〕\jruby[g]{獨得}{どくとく}の\jruby[g]{見}{けん}は\jruby[g]{私}{わたくし}に似る、人其の\jruby[g]{驟至}{しうし}に\jruby[g]{驚}{おどろ}く。\jruby[g]{平凡}{へいぼん}の\jruby[g]{議}{ぎ}は公に似る、世其の\jruby[g]{狃聞}{ぢうぶん}に安んず。凡そ人の言を\jruby[g]{聽}{き}くは、宜しく\jruby[g]{虚懷}{きよくわい}にして之を\jruby[g]{邀}{むか}ふべし。\jruby[g]{狃聞}{ぢうぶん}に\jruby[g]{苟安}{こうあん}することなくんば可なり。

\end{azconvIndented}

\begin{azconvIndented}{1\zw}{-1\zw}

六五　心理是豎工夫、愽覽是横工夫。豎工夫、則深入自得。横工夫、則淺易汎濫。

\end{azconvIndented}
\begin{azconvIndented}{2\zw}{-1\zw}

〔譯〕\jruby[g]{心理}{しんり}は是れ\jruby[g]{豎}{たて}の工夫なり、\jruby[g]{愽覽}{はくらん}は是れ\jruby[g]{横}{よこ}の工夫なり。\jruby[g]{豎}{たて}の工夫は、則ち\jruby[g]{深入}{しんにふ}\jruby[g]{自得}{じとく}せよ。\jruby[g]{横}{よこ}の工夫は、則ち\jruby[g]{淺易}{せんい}\jruby[g]{汎濫}{はんらん}なれ。

\end{azconvIndented}

\begin{azconvIndented}{1\zw}{-1\zw}

六六　讀\azconvKunten{レ}經、宜\azconvKunten{下}以\azconvKunten{二}我之心\azconvKunten{一}讀\azconvKunten{二}經之心\azconvKunten{一}、以\azconvKunten{二}經之心\azconvKunten{一}釋\azconvKunten{中}我之心\azconvKunten{上}。不\azconvKunten{レ}然徒爾講\azconvKunten{二}明訓詁\azconvKunten{一}而已、便是終身不\azconvKunten{二}曾讀\azconvKunten{一}。

\end{azconvIndented}
\begin{azconvIndented}{2\zw}{-1\zw}

〔譯〕\jruby[g]{經}{けい}を讀むは、宜しく我れの心を以て經の心を讀み、經の心を以て我の心を\jruby[g]{釋}{しやく}すべし。然らずして\jruby[g]{徒爾}{とじ}に\jruby[g]{訓詁}{くんこ}を\jruby[g]{講明}{かうめい}するのみならば、\jruby[g]{便}{すなは}ち是れ終身\jruby[g]{曾}{かつ}て讀まざるなり。

\end{azconvIndented}

\begin{azconvIndented}{1\zw}{-1\zw}

六七　引\azconvKunten{レ}滿中\azconvKunten{レ}度、發無\azconvKunten{二}空箭\azconvKunten{一}。人事宜\azconvKunten{二}如\azconvKunten{レ}射然\azconvKunten{一}。

\end{azconvIndented}
\begin{azconvIndented}{2\zw}{-1\zw}

〔譯〕\jruby[g]{滿}{まん}を\jruby[g]{引}{ひ}き\jruby[g]{度}{ど}に\jruby[g]{中}{あた}り、發して\jruby[g]{空箭}{くうぜん}無し。人事宜しく\jruby[g]{射}{しや}の如く然るべし。

\end{azconvIndented}

\begin{azconvIndented}{1\zw}{-1\zw}

六八　前人、謂\azconvKunten{二}英氣害\azconvKunten{\azconvIR}事。余則謂、英氣不\azconvKunten{レ}可\azconvKunten{レ}無、但露\azconvKunten{二}圭角\azconvKunten{一}爲\azconvKunten{二}不可\azconvKunten{一}。

\end{azconvIndented}
\begin{azconvIndented}{2\zw}{-1\zw}

譯前人は、\jruby[g]{英氣}{えいき}は事を\jruby[g]{害}{がい}すと謂へり。余は則ち謂ふ、英氣は無かる可らず、\jruby[g]{但}{た}だ\jruby[g]{圭角}{けいかく}を\jruby[g]{露}{あら}はすを不可と爲すと。

\end{azconvIndented}

\begin{azconvIndented}{1\zw}{-1\zw}

六九　刀槊之技、懷\azconvKunten{二}怯心\azconvKunten{一}者衄、頼\azconvKunten{二}勇氣\azconvKunten{一}者敗。必也泯\azconvKunten{二}勇怯於一靜\azconvKunten{一}、忘\azconvKunten{二}勝負於一動\azconvKunten{一}。動\azconvKunten{レ}之以\azconvKunten{レ}天、廓然太公、靜\azconvKunten{レ}之以\azconvKunten{レ}地、物來順應。如\azconvKunten{レ}是者勝矣。心學亦不\azconvKunten{レ}外\azconvKunten{二}於此\azconvKunten{一}。

\end{azconvIndented}
\begin{azconvIndented}{2\zw}{-1\zw}

〔譯〕\jruby[g]{刀槊}{たうさく}の\jruby[g]{技}{ぎ}、\jruby[g]{怯}{きよ}心を\jruby[g]{懷}{いだ}く者は\jruby[g]{衄}{くじ}け、\jruby[g]{勇氣}{ゆうき}を\jruby[g]{頼}{たの}む者は\jruby[g]{敗}{やぶ}る。必や\jruby[g]{勇怯}{ゆうきよ}を一\jruby[g]{靜}{せい}に\jruby[g]{泯}{ほろぼ}し、\jruby[g]{勝負}{しようぶ}を一\jruby[g]{動}{どう}に\jruby[g]{忘}{わす}れ、之を\jruby[g]{動}{うご}かすに天を以てして、\jruby[g]{廓然}{かくぜん}\jruby[g]{太公}{たいこう}に、之を\jruby[g]{靜}{しづ}むるに地を以てして、\jruby[g]{物}{もの}來つて\jruby[g]{順應}{じゆんおう}せん。\jruby[g]{是}{かく}の如き者は\jruby[g]{勝}{か}たん。心學も亦\jruby[g]{此}{こゝ}に外ならず。

\end{azconvIndented}
\begin{azconvIndented}{3\zw}{-1\zw}

評長兵京師に\jruby[g]{敗}{やぶ}る。木戸公は岡部氏に\jruby[g]{寄}{よ}つて\jruby[g]{禍}{わざはい}を\jruby[g]{免}{まぬか}るゝことを得たり。\jruby[g]{後}{のち}丹波に\jruby[g]{赴}{おもむ}き、\jruby[g]{姓名}{せいめい}を\jruby[g]{變}{か}へ、\jruby[g]{博徒}{ばくと}に\jruby[g]{混}{まじ}り、\jruby[g]{酒客}{しゆかく}に\jruby[g]{交}{まじは}り、以て時勢を\jruby[g]{窺}{うかゞ}へり。南洲は\jruby[g]{浪華}{なには}の某樓に\jruby[g]{寓}{ぐう}す。幕吏\jruby[g]{搜索}{さうさく}して樓下に至る。南洲乃ち\jruby[g]{劇}{げき}を觀るに託して、舟を\jruby[g]{僦}{か}りて\jruby[g]{逃}{に}げ去れり。此れ皆\jruby[g]{勇怯}{ゆうきよ}を\jruby[g]{泯}{ほろぼ}し\jruby[g]{勝負}{しようぶ}を忘るゝものなり。

\end{azconvIndented}

\begin{azconvIndented}{1\zw}{-1\zw}

七〇　無\azconvKunten{レ}我則不\azconvKunten{レ}獲\azconvKunten{二}其身\azconvKunten{一}、即是義。無\azconvKunten{レ}物則不\azconvKunten{レ}見\azconvKunten{二}其人\azconvKunten{一}、即是勇。

\end{azconvIndented}
\begin{azconvIndented}{2\zw}{-1\zw}

〔譯〕\jruby[g]{我}{わ}れ無ければ則ち其身を\jruby[g]{獲}{え}ず、即ち是れ\jruby[g]{義}{ぎ}なり。物無ければ則ち其人を見ず、即ち是れ\jruby[g]{勇}{ゆう}なり。

\end{azconvIndented}

\begin{azconvIndented}{1\zw}{-1\zw}

七一　自反而縮者、無\azconvKunten{レ}我也。雖\azconvKunten{二}千萬人\azconvKunten{一}吾往矣、無\azconvKunten{レ}物也。

\end{azconvIndented}
\begin{azconvIndented}{2\zw}{-1\zw}

譯自ら\jruby[g]{反}{かへり}みて\jruby[g]{縮}{なほ}きは、\jruby[g]{我}{われ}無きなり。千萬人と雖吾れ往かんは、物無きなり。

\end{azconvIndented}

\begin{azconvIndented}{1\zw}{-1\zw}

七二　三軍不\azconvKunten{レ}和、難\azconvKunten{二}以言\azconvKunten{\azconvIR}戰。百官不\azconvKunten{レ}和、難\azconvKunten{二}以言\azconvKunten{\azconvIR}治。書云、同\azconvKunten{レ}寅協\azconvKunten{レ}恭和衷哉。唯一和字、一\azconvKunten{二}串治亂\azconvKunten{一}。

\end{azconvIndented}
\begin{azconvIndented}{2\zw}{-1\zw}

譯三軍和せずば、以て\jruby[g]{戰}{たゝかひ}を言ひ\jruby[g]{難}{がた}し。百官和せずば、以て\jruby[g]{治}{ち}を言ひ難し。書に云ふ、\jruby[g]{寅}{いん}を同じうし\jruby[g]{恭}{きよう}を\jruby[g]{協}{あは}せ\jruby[g]{和衷}{わちゆう}せよやと。唯だ一の和字、\jruby[g]{治亂}{ちらん}を\jruby[g]{一串}{いつくわん}す。

\end{azconvIndented}
\begin{azconvIndented}{3\zw}{-1\zw}

〔評〕\jruby[g]{復古}{ふくこ}の\jruby[g]{業}{げふ}は\jruby[g]{薩長}{さつちやう}の\jruby[g]{合縱}{がつしよう}に成る。是れより先き、土人坂本\jruby[g]{龍馬}{りゆうま}、薩長の和せざるを\jruby[g]{憂}{うれ}へ、薩\jruby[g]{邸}{てい}に\jruby[g]{抵}{いた}り、大久保・西郷諸氏に説き、又長邸に\jruby[g]{抵}{いた}り、木戸・大村諸氏に説く。薩人黒田・大山諸氏長に至り、長人木戸・品川諸氏薩に\jruby[g]{往}{ゆ}き、而て後\jruby[g]{和}{わ}成り、\jruby[g]{維新}{いしん}の\jruby[g]{鴻業}{こうげふ}を\jruby[g]{致}{いた}せり。

\end{azconvIndented}

\begin{azconvIndented}{1\zw}{-1\zw}

七三　凡事有\azconvKunten{二}眞是非\azconvKunten{一}、有\azconvKunten{二}假是非\azconvKunten{一}。假是非、謂\azconvKunten{三}通俗之所\azconvKunten{二}可否\azconvKunten{一}。年少未\azconvKunten{レ}學、而先了\azconvKunten{二}假是非\azconvKunten{一}、迨\azconvKunten{レ}後欲\azconvKunten{レ}得\azconvKunten{二}眞是非\azconvKunten{一}、亦不\azconvKunten{レ}易\azconvKunten{レ}入。所\azconvKunten{レ}謂先入爲\azconvKunten{レ}主、不\azconvKunten{レ}可\azconvKunten{二}如何\azconvKunten{一}耳。

\end{azconvIndented}
\begin{azconvIndented}{2\zw}{-1\zw}

譯凡そ事に\jruby[g]{眞是非}{しんぜひ}有り、\jruby[g]{假是非}{かぜひ}有り。假是非とは、\jruby[g]{通俗}{つうぞく}の可否する所を謂ふ。年\jruby[g]{少}{わか}く未だ學ばずして、先づ假是非を\jruby[g]{了}{れう}し、後に\jruby[g]{迨}{およ}んで眞是非を得んと欲するも、亦入り\jruby[g]{易}{やす}からず。謂はゆる\jruby[g]{先入}{せんにふ}\jruby[g]{主}{しゆ}と\jruby[g]{爲}{な}り、如何ともす可らざるのみ。

\end{azconvIndented}

\begin{azconvIndented}{1\zw}{-1\zw}

七四　果斷、有\azconvKunten{二}自\azconvKunten{レ}義來者\azconvKunten{一}。有\azconvKunten{二}自\azconvKunten{レ}智來者\azconvKunten{一}。有\azconvKunten{二}自\azconvKunten{レ}勇來者\azconvKunten{一}。有\azconvKunten{下}并\azconvKunten{二}義與\azconvKunten{\azconvIR}智而來者\azconvKunten{上}、上也。徒勇而已者殆矣。

\end{azconvIndented}
\begin{azconvIndented}{2\zw}{-1\zw}

〔譯〕\jruby[g]{果斷}{くわだん}は、\jruby[g]{義}{ぎ}より來るもの有り。\jruby[g]{智}{ち}より來るもの有り。\jruby[g]{勇}{ゆう}より來るもの有り。義と智とを\jruby[g]{併}{あは}せて來るもの有り、\jruby[g]{上}{じやう}なり。\jruby[g]{徒}{たゞ}に\jruby[g]{勇}{ゆう}のみなるは\jruby[g]{殆}{あやふ}し。

\end{azconvIndented}
\begin{azconvIndented}{3\zw}{-1\zw}

〔評〕\jruby[g]{關八州}{くわんはつしう}は古より武を用ふるの地と稱す。\jruby[g]{興世}{おきよ}王\jruby[g]{反逆}{はんぎやく}すと雖、猶\jruby[g]{將門}{まさかど}に説いて之に\jruby[g]{據}{よ}らしむ。小田原の\jruby[g]{役}{えき}、\jruby[g]{豐}{ほう}公は徳川公に謂うて曰ふ、東方に地あり、\jruby[g]{江戸}{えど}と曰ふ、以て\jruby[g]{都府}{とふ}を開く可しと。\jruby[g]{一新}{いつしん}の\jruby[g]{始}{はじ}め、大久保公\jruby[g]{遷都}{せんと}の\jruby[g]{議}{ぎ}を\jruby[g]{獻}{けん}じて曰ふ、官軍已に\jruby[g]{勝}{か}つと雖、\jruby[g]{東賊}{とうぞく}猶未だ\jruby[g]{滅}{ほろ}びず、宜しく\jruby[g]{非常}{ひじやう}の\jruby[g]{斷}{だん}を以て非常の事を行ふべしと。先見の明\jruby[g]{智}{ち}と謂ふ可し。

\end{azconvIndented}

\begin{azconvIndented}{1\zw}{-1\zw}

七五　公私在\azconvKunten{レ}事、又在\azconvKunten{レ}情。事公而情私者有\azconvKunten{レ}之。事私而情公者有\azconvKunten{レ}之。爲\azconvKunten{レ}政者、宜\azconvKunten{下}權\azconvKunten{二}衡人情事理輕重處\azconvKunten{一}、以用\azconvKunten{中}其中於\azconvKunten{\azconvJR}民。

\end{azconvIndented}
\begin{azconvIndented}{2\zw}{-1\zw}

〔譯〕\jruby[g]{公私}{こうし}は事に在り、又情に在り。事公にして情私なるもの之有り。事私にして情公なるもの之有り。政を爲す者は、宜しく人情\jruby[g]{事理}{じり}\jruby[g]{輕重}{けいぢゆう}の處を\jruby[g]{權衡}{けんかう}して、以て其の\jruby[g]{中}{ちゆう}を民に用ふべし。

\end{azconvIndented}
\begin{azconvIndented}{3\zw}{-1\zw}

評南洲城山に\jruby[g]{據}{よ}る。官軍\jruby[g]{柵}{さく}を\jruby[g]{植}{う}ゑて之を守る。\jruby[g]{山縣}{やまがた}中將書を南洲に寄せて兩軍\jruby[g]{殺傷}{さつしやう}の\jruby[g]{慘}{さん}を\jruby[g]{極言}{きよくげん}す。南洲其の書を見て曰ふ、我れ山縣に\jruby[g]{負}{そむ}かずと、\jruby[g]{斷然}{だんぜん}死に\jruby[g]{就}{つ}けり。中將は南洲の\jruby[g]{元}{げん}を\jruby[g]{視}{み}て曰ふ、\jruby[g]{惜}{を}しいかな、天下の一勇將を失へりと、\jruby[g]{流涕}{りうてい}すること之を久しうせり。\jruby[g]{噫}{あゝ}公私情盡せり。

\end{azconvIndented}

\begin{azconvIndented}{1\zw}{-1\zw}

七六　愼獨工夫、當\azconvKunten{下}如\azconvKunten{三}身在\azconvKunten{二}稠人廣座中\azconvKunten{一}一般\azconvKunten{上}。應酬工夫、當\azconvKunten{下}如\azconvKunten{二}間居獨處時\azconvKunten{一}一般\azconvKunten{上}。

\end{azconvIndented}
\begin{azconvIndented}{2\zw}{-1\zw}

〔譯〕\jruby[g]{愼獨}{しんどく}の\jruby[g]{工夫}{くふう}は、\jruby[g]{當}{まさ}に身\jruby[g]{稠人}{ちうじん}\jruby[g]{廣座}{くわうざ}の中に在るが如く一\jruby[g]{般}{ぱん}なるべし。\jruby[g]{應酬}{おうしう}の工夫は、\jruby[g]{當}{まさ}に\jruby[g]{間居}{かんきよ}\jruby[g]{獨處}{どくしよ}の時の如く一般なるべし。

\end{azconvIndented}

\begin{azconvIndented}{1\zw}{-1\zw}

七七　心要\azconvKunten{二}現在\azconvKunten{一}。事未\azconvKunten{レ}來、不\azconvKunten{レ}可\azconvKunten{レ}邀。事已往、不\azconvKunten{レ}可\azconvKunten{レ}追。纔追纔邀、便是放心。

\end{azconvIndented}
\begin{azconvIndented}{2\zw}{-1\zw}

譯心は\jruby[g]{現在}{げんざい}せんことを\jruby[g]{要}{えう}す。事未だ來らずば、\jruby[g]{邀}{むか}ふ可らず。事已に\jruby[g]{往}{ゆ}かば、\jruby[g]{追}{お}ふ可らず。\jruby[g]{纔}{わづ}かに追ひ纔かに邀へば、\jruby[g]{便}{すなは}ち是れ\jruby[g]{放心}{はうしん}なり。

\end{azconvIndented}

\begin{azconvIndented}{1\zw}{-1\zw}

七八　物集\azconvKunten{二}於其所\azconvKunten{\azconvIR}好、人也。事赴\azconvKunten{二}於所\azconvKunten{\azconvIR}不\azconvKunten{レ}期、天也。

\end{azconvIndented}
\begin{azconvIndented}{2\zw}{-1\zw}

〔譯〕\jruby[g]{物}{もの}其の好む所に\jruby[g]{集}{あつま}るは、人なり。\jruby[g]{事}{こと}\jruby[g]{期}{き}せざる所に\jruby[g]{赴}{おもむ}くは、天なり。

\end{azconvIndented}

\begin{azconvIndented}{1\zw}{-1\zw}

七九　人貴\azconvKunten{二}厚重\azconvKunten{一}、不\azconvKunten{レ}貴\azconvKunten{二}遲重\azconvKunten{一}。尚\azconvKunten{二}眞率\azconvKunten{一}、不\azconvKunten{レ}尚\azconvKunten{二}輕率\azconvKunten{一}。

\end{azconvIndented}
\begin{azconvIndented}{2\zw}{-1\zw}

譯人は、\jruby[g]{厚重}{こうちよう}を貴ぶ、\jruby[g]{遲重}{ちちよう}を貴ばず。\jruby[g]{眞率}{しんそつ}を\jruby[g]{尚}{たつと}ぶ、\jruby[g]{輕率}{けいそつ}を尚ばず。

\end{azconvIndented}
\begin{azconvIndented}{3\zw}{-1\zw}

評南洲人に\jruby[g]{接}{せつ}して、\jruby[g]{妄}{みだり}に\jruby[g]{語}{ご}を\jruby[g]{交}{まじ}へず、人之を\jruby[g]{憚}{はゞか}る。然れども其の人を知るに及んでは、則ち心を\jruby[g]{傾}{かたむ}けて之を\jruby[g]{援}{たす}く。其人に非ざれば則ち\jruby[g]{終身}{しゆうしん}\jruby[g]{言}{い}はず。

\end{azconvIndented}

\begin{azconvIndented}{1\zw}{-1\zw}

八〇　凡生物皆資\azconvKunten{二}於養\azconvKunten{一}。天生而地養\azconvKunten{レ}之。人則地之氣精英。吾欲\azconvKunten{三}靜坐以養\azconvKunten{レ}氣、動行以養\azconvKunten{レ}體、氣體相資、以養\azconvKunten{二}此生\azconvKunten{一}。所\azconvKunten{二}以從\azconvKunten{レ}地而事\azconvKunten{\azconvIR}天。

\end{azconvIndented}
\begin{azconvIndented}{2\zw}{-1\zw}

譯凡そ生物は皆\jruby[g]{養}{やう}を\jruby[g]{資}{と}る。天生じて地之を\jruby[g]{養}{やしな}ふ。人は則ち地の氣の\jruby[g]{精英}{せいえい}なり。吾れ靜坐して以て氣を養ひ、\jruby[g]{動行}{どうかう}して以て體を養ひ、氣と體と相\jruby[g]{資}{と}つて以て此の生を養はんと欲す。地に從うて天に事ふる所以なり。

\end{azconvIndented}
\begin{azconvIndented}{3\zw}{-1\zw}

評維新の\jruby[g]{業}{げふ}は三藩の兵力に由ると雖、抑之を養ふに\jruby[g]{素}{そ}あり、曰く\jruby[g]{名義}{めいぎ}なり、曰く\jruby[g]{名分}{めいぶん}なり。或は云ふ、維新の\jruby[g]{功}{こう}は\jruby[g]{大日本史}{だいにつぽんし}及び外史に\jruby[g]{基}{もと}づくと、亦\jruby[g]{理}{り}\jruby[g]{無}{な}しとせざるなり。

\end{azconvIndented}

\begin{azconvIndented}{1\zw}{-1\zw}

八一　凡爲\azconvKunten{レ}學之初、必立\azconvKunten{下}欲\azconvKunten{レ}爲\azconvKunten{二}大人\azconvKunten{一}之志\azconvKunten{上}、然後書可\azconvKunten{レ}讀也。不\azconvKunten{レ}然、徒貪\azconvKunten{二}聞見\azconvKunten{一}而已、則或恐\azconvKunten{二}長\azconvKunten{レ}傲飾\azconvKunten{\azconvIR}非。所\azconvKunten{レ}謂假\azconvKunten{二}寇兵\azconvKunten{一}、資\azconvKunten{二}盜糧\azconvKunten{一}也、可\azconvKunten{レ}虞。

\end{azconvIndented}
\begin{azconvIndented}{2\zw}{-1\zw}

譯凡そ學を爲すの初め、必ず大人たらんと欲するの志を立て、然る後書讀む可し。然らずして、\jruby[g]{徒}{いたづら}に聞見を\jruby[g]{貪}{むさぼ}るのみならば、則ち或は\jruby[g]{傲}{がう}を\jruby[g]{長}{ちやう}じ非を\jruby[g]{飾}{かざ}らんことを恐る。謂はゆる\jruby[g]{寇}{こう}に兵を\jruby[g]{假}{か}し、\jruby[g]{盜}{たう}に\jruby[g]{糧}{りやう}を\jruby[g]{資}{し}するなり、\jruby[g]{虞}{おもんぱか}る可し。

\end{azconvIndented}

\begin{azconvIndented}{1\zw}{-1\zw}

八二　以\azconvKunten{二}眞己\azconvKunten{一}克\azconvKunten{二}假己\azconvKunten{一}、天理也。以\azconvKunten{二}身我\azconvKunten{一}害\azconvKunten{二}心我\azconvKunten{一}、人欲也。

\end{azconvIndented}
\begin{azconvIndented}{2\zw}{-1\zw}

〔譯〕\jruby[g]{眞己}{しんこ}を以て\jruby[g]{假己}{かこ}に\jruby[g]{克}{か}つ、天理なり。\jruby[g]{身我}{しんが}を以て心我を\jruby[g]{害}{がい}す、\jruby[g]{人欲}{じんよく}なり。

\end{azconvIndented}

\begin{azconvIndented}{1\zw}{-1\zw}

八三　無\azconvKunten{二}一息間斷\azconvKunten{一}、無\azconvKunten{二}一刻急忙\azconvKunten{一}。即是天地氣象。

\end{azconvIndented}
\begin{azconvIndented}{2\zw}{-1\zw}

譯一\jruby[g]{息}{そく}の\jruby[g]{間斷}{かんだん}無く、一\jruby[g]{刻}{こく}の\jruby[g]{急忙}{きふばう}無し。即ち是れ天地の\jruby[g]{氣象}{きしやう}なり。

\end{azconvIndented}
\begin{azconvIndented}{3\zw}{-1\zw}

評木戸公毎旦\jruby[g]{考妣}{ちゝはゝ}の木主を拜す。身\jruby[g]{煩劇}{はんげき}に居ると雖、少しくも\jruby[g]{怠}{おこた}らず。三十年の間一日の如し。

\end{azconvIndented}

\begin{azconvIndented}{1\zw}{-1\zw}

八四　有\azconvKunten{レ}心\azconvKunten{二}於無\azconvKunten{\azconvIR}心、工夫是也。無\azconvKunten{レ}心\azconvKunten{二}於有\azconvKunten{\azconvIR}心、本體是也。

\end{azconvIndented}
\begin{azconvIndented}{2\zw}{-1\zw}

譯心無きに心有るは、\jruby[g]{工夫}{くふう}是なり。心有るに心無きは、\jruby[g]{本體}{ほんたい}是なり。

\end{azconvIndented}

\begin{azconvIndented}{1\zw}{-1\zw}

八五　不\azconvKunten{レ}知而知者、道心也。知而不\azconvKunten{レ}知者、人心也。

\end{azconvIndented}
\begin{azconvIndented}{2\zw}{-1\zw}

譯知らずして知る者は、\jruby[g]{道心}{だうしん}なり。知つて知らざる者は、\jruby[g]{人心}{じんしん}なり。

\end{azconvIndented}

\begin{azconvIndented}{1\zw}{-1\zw}

八六　心靜、方能知\azconvKunten{二}白日\azconvKunten{一}。眼明、始會\azconvKunten{レ}識\azconvKunten{二}青天\azconvKunten{一}。此程伯氏之句也。青天白日、常在\azconvKunten{二}於我\azconvKunten{一}。宜\azconvKunten{下}掲\azconvKunten{二}之座右\azconvKunten{一}、以爲\azconvKunten{中}警戒\azconvKunten{上}。

\end{azconvIndented}
\begin{azconvIndented}{2\zw}{-1\zw}

譯心\jruby[g]{靜}{しづか}にして、\jruby[g]{方}{まさ}に能く白日を知る。眼明かにして、始めて青天を識り\jruby[g]{會}{え}すと。此れ\jruby[g]{程伯氏}{ていはくし}の句なり。青天白日は、常に我に在り。宜しく之を\jruby[g]{座右}{ざいう}に\jruby[g]{掲}{かゝ}げて、以て\jruby[g]{警戒}{けいかい}と爲すべし。

\end{azconvIndented}

\begin{azconvIndented}{1\zw}{-1\zw}

八七　靈光充\azconvKunten{レ}體時、細大事物、無\azconvKunten{二}遺落\azconvKunten{一}、無\azconvKunten{二}遲疑\azconvKunten{一}。

\end{azconvIndented}
\begin{azconvIndented}{2\zw}{-1\zw}

〔譯〕\jruby[g]{靈光}{れいくわう}\jruby[g]{體}{たい}に\jruby[g]{充}{み}つる時、\jruby[g]{細大}{さいだい}の事物、\jruby[g]{遺落}{ゐらく}無く、\jruby[g]{遲疑}{ちぎ}無し。

\end{azconvIndented}
\begin{azconvIndented}{3\zw}{-1\zw}

評死を決するは、\jruby[g]{薩}{さつ}の長ずる所なり。公義を説くは、土の\jruby[g]{俗}{ぞく}なり。\jruby[g]{維新}{いしん}の初め、一公卿あり、南洲の所に往いて\jruby[g]{復古}{ふくこ}の事を説く。南洲曰ふ、夫れ復古は\jruby[g]{易事}{いじ}に非ず、且つ九重\jruby[g]{阻絶}{そぜつ}し、\jruby[g]{妄}{みだり}に藩人を通ずるを得ず、必ずや\jruby[g]{縉紳}{しんしん}死を致す有らば、則ち事或は成らんと。又\jruby[g]{後藤象}{ごとうしやう}次郎に\jruby[g]{往}{ゆ}いて之を説く。象次郎曰ふ、復古は\jruby[g]{難}{かた}きに非ず、然れども\jruby[g]{門地}{もんち}を\jruby[g]{廢}{はい}し、\jruby[g]{門閥}{もんばつ}を\jruby[g]{罷}{や}め、\jruby[g]{賢}{けん}を\jruby[g]{擧}{あ}ぐること\jruby[g]{方}{はう}なきに非ざれば、則ち不可なりと。二人の本領自ら\jruby[g]{見}{あら}はる。

\end{azconvIndented}

\begin{azconvIndented}{1\zw}{-1\zw}

八八　人心之靈、如\azconvKunten{二}太陽\azconvKunten{一}然。但克伐怨欲、雲霧四塞、此靈烏在。故誠\azconvKunten{レ}意工夫、莫\azconvKunten{レ}先\azconvKunten{下}於掃\azconvKunten{二}雲霧\azconvKunten{一}仰\azconvKunten{中}白日\azconvKunten{上}。凡爲\azconvKunten{レ}學之要、自\azconvKunten{レ}此而起\azconvKunten{レ}基。故曰、誠者物之終始。

\end{azconvIndented}
\begin{azconvIndented}{2\zw}{-1\zw}

譯人心の\jruby[g]{靈}{れい}、\jruby[g]{太陽}{たいやう}の如く然り。但だ\jruby[g]{克伐}{こくばつ}\jruby[g]{怨欲}{えんよく}、\jruby[g]{雲霧}{うんむ}\jruby[g]{四塞}{しそく}せば、此の\jruby[g]{靈}{れい}\jruby[g]{烏}{いづ}くに在る。故に意を\jruby[g]{誠}{まこと}にする工夫は、\jruby[g]{雲霧}{うんむ}を\jruby[g]{掃}{はら}うて白日を\jruby[g]{仰}{あふ}ぐより先きなるは\jruby[g]{莫}{な}し。凡そ學を爲すの\jruby[g]{要}{えう}は、\jruby[g]{此}{これ}よりして\jruby[g]{基}{もとゐ}を\jruby[g]{起}{おこ}す。故に曰ふ、誠は物の\jruby[g]{終始}{しゆうし}と。

\end{azconvIndented}

\begin{azconvIndented}{1\zw}{-1\zw}

八九　胸次清快、則人事百艱亦不\azconvKunten{レ}阻。

\end{azconvIndented}
\begin{azconvIndented}{2\zw}{-1\zw}

〔譯〕\jruby[g]{胸次}{きようじ}\jruby[g]{清快}{せいくわい}なれば、則ち人事百\jruby[g]{艱}{かん}亦\jruby[g]{阻}{そ}せず。

\end{azconvIndented}

\begin{azconvIndented}{1\zw}{-1\zw}

九〇　人心之靈、主\azconvKunten{二}於氣\azconvKunten{一}。氣體之充也。凡爲\azconvKunten{レ}事、以\azconvKunten{レ}氣爲\azconvKunten{二}先導\azconvKunten{一}、則擧體無\azconvKunten{二}失措\azconvKunten{一}。技能工藝、亦皆如\azconvKunten{レ}此。

\end{azconvIndented}
\begin{azconvIndented}{2\zw}{-1\zw}

譯人心の\jruby[g]{靈}{れい}は、\jruby[g]{氣}{き}を\jruby[g]{主}{しゆ}とす。氣は\jruby[g]{體}{たい}に之れ\jruby[g]{充}{み}つるものなり。凡そ事を爲すに、氣を以て\jruby[g]{先導}{せんだう}と爲さば、則ち\jruby[g]{擧體}{きよたい}\jruby[g]{失措}{しつそ}無し。\jruby[g]{技能}{ぎのう}\jruby[g]{工藝}{こうげい}も、亦皆\jruby[g]{此}{かく}の如し。

\end{azconvIndented}

\begin{azconvIndented}{1\zw}{-1\zw}

九一　靈光無\azconvKunten{二}障碍\azconvKunten{一}、則氣乃流動不\azconvKunten{レ}餒、四體覺\azconvKunten{レ}輕。

\end{azconvIndented}
\begin{azconvIndented}{2\zw}{-1\zw}

〔譯〕\jruby[g]{靈光}{れいくわう}\jruby[g]{障碍}{しやうげ}無くば、則ち\jruby[g]{氣}{き}乃ち\jruby[g]{流動}{りうどう}して\jruby[g]{餒}{う}ゑず、\jruby[g]{四體}{したい}\jruby[g]{輕}{かる}きを\jruby[g]{覺}{おぼ}えん。

\end{azconvIndented}

\begin{azconvIndented}{1\zw}{-1\zw}

九二　英氣是天地精英之氣。聖人薀\azconvKunten{二}之於内\azconvKunten{一}、不\azconvKunten{三}肯露\azconvKunten{二}諸外\azconvKunten{一}。賢者則時時露\azconvKunten{レ}之。自餘豪傑之士、全然露\azconvKunten{レ}之。若\azconvKunten{下}夫絶無\azconvKunten{二}此氣\azconvKunten{一}者\azconvKunten{上}、爲\azconvKunten{二}鄙夫小人\azconvKunten{一}、碌碌不\azconvKunten{レ}足\azconvKunten{レ}算者爾。

\end{azconvIndented}
\begin{azconvIndented}{2\zw}{-1\zw}

譯英氣は是れ天地\jruby[g]{精英}{せいえい}の氣なり。聖人は之を内に\jruby[g]{薀}{をさ}めて、\jruby[g]{肯}{あへ}て\jruby[g]{諸}{これ}を外に\jruby[g]{露}{あら}はさず。賢者は則ち時時之を\jruby[g]{露}{あら}はす。\jruby[g]{自餘}{じよ}豪傑の士は、全然之を\jruby[g]{露}{あら}はす。\jruby[g]{夫}{か}の\jruby[g]{絶}{た}えて\jruby[g]{此}{この}\jruby[g]{氣}{き}なき者の若きは、\jruby[g]{鄙夫}{ひふ}小人と爲す、\jruby[g]{碌碌}{ろく{〳〵}}として\jruby[g]{算}{かぞ}ふるに足らざるもののみ。

\end{azconvIndented}

\begin{azconvIndented}{1\zw}{-1\zw}

九三　人須\azconvKunten{レ}著\azconvKunten{二}忙裏占\azconvKunten{レ}間、苦中存\azconvKunten{レ}樂工夫\azconvKunten{一}。

\end{azconvIndented}
\begin{azconvIndented}{2\zw}{-1\zw}

譯人は須らく\jruby[g]{忙裏}{ばうり}に\jruby[g]{間}{かん}を\jruby[g]{占}{し}め、\jruby[g]{苦中}{くちゆう}に\jruby[g]{樂}{らく}を存ずる工夫を\jruby[g]{著}{つ}くべし。

\end{azconvIndented}
\begin{azconvIndented}{3\zw}{-1\zw}

評南洲岩崎谷洞中に居る。砲丸雨の如く、洞口を出づる能はず。詩あり云ふ「百戰無\azconvKunten{レ}功半歳間、首邱幸得\azconvKunten{レ}返\azconvKunten{二}家山\azconvKunten{一}。笑儂向\azconvKunten{レ}死如\azconvKunten{二}仙客\azconvKunten{一}。盡日洞中棋響間」（編者曰、此詩、長州ノ人杉孫七郎ノ作ナリ、南洲翁ノ作ト稱スルハ誤ル）謂はゆる\jruby[g]{忙}{ばう}中に間を占むる者なり。然れども亦以て其の戰志無きを知るべし。余句あり、云ふ「可\azconvKunten{レ}見南洲無\azconvKunten{二}戰志\azconvKunten{一}。砲丸雨裡間牽\azconvKunten{レ}犬」と、是れ\jruby[g]{實録}{じつろく}なり。

\end{azconvIndented}

\begin{azconvIndented}{1\zw}{-1\zw}

九四　凡區\azconvKunten{二}處人事\azconvKunten{一}、當\azconvKunten{下}先慮\azconvKunten{二}其結局處\azconvKunten{一}、而後下\azconvKunten{\azconvJR}手。無\azconvKunten{レ}楫之舟勿\azconvKunten{レ}行、無\azconvKunten{レ}的之箭勿\azconvKunten{レ}發。

\end{azconvIndented}
\begin{azconvIndented}{2\zw}{-1\zw}

譯凡そ人事を\jruby[g]{區處}{くしよ}するには、當さに先づ其の\jruby[g]{結局}{けつきよく}の處を\jruby[g]{慮}{おもんぱ}かりて、後に手を下すべし。\jruby[g]{楫}{かぢ}無きの舟は\jruby[g]{行}{や}る\jruby[g]{勿}{なか}れ、\jruby[g]{的}{まと}無きの\jruby[g]{箭}{や}は\jruby[g]{發}{はな}つ勿れ。

\end{azconvIndented}

\begin{azconvIndented}{1\zw}{-1\zw}

九五　朝而不\azconvKunten{レ}食、則晝而饑。少而不\azconvKunten{レ}學、則壯而惑。饑者猶可\azconvKunten{レ}忍、惑者不\azconvKunten{レ}可\azconvKunten{二}奈何\azconvKunten{一}。

\end{azconvIndented}
\begin{azconvIndented}{2\zw}{-1\zw}


〔譯〕\jruby[g]{朝}{あさ}にして\jruby[g]{食}{くら}はずば、\jruby[g]{晝}{ひる}にして\jruby[g]{饑}{う}う。\jruby[g]{少}{わか}うして學ばずば、壯にして\jruby[g]{惑}{まど}ふ。饑うるは猶\jruby[g]{忍}{しの}ぶ可し、\jruby[g]{惑}{まど}ふは奈何ともす可からず。

\end{azconvIndented}

\begin{azconvIndented}{1\zw}{-1\zw}

九六　今日之貧賤不\azconvKunten{レ}能\azconvKunten{二}素行\azconvKunten{一}、乃他日之富貴、必驕泰。今日之富貴不\azconvKunten{レ}能\azconvKunten{二}素行\azconvKunten{一}、乃他日之患難、必狼狽。

\end{azconvIndented}
\begin{azconvIndented}{2\zw}{-1\zw}

譯今日の\jruby[g]{貧賤}{ひんせん}に\jruby[g]{素行}{そかう}する能はずば、乃ち他日の\jruby[g]{富貴}{ふうき}に、必ず\jruby[g]{驕泰}{けうたい}ならん。今日の\jruby[g]{富貴}{ふうき}に\jruby[g]{素行}{そかう}する能はずんば、乃ち他日の\jruby[g]{患難}{くわんなん}に、必ず\jruby[g]{狼狽}{らうばい}せん。

\end{azconvIndented}
\begin{azconvIndented}{3\zw}{-1\zw}

評南洲、\jruby[g]{顯職}{けんしよく}に居り\jruby[g]{勳功}{くんこう}を\jruby[g]{負}{お}ふと雖、身極めて\jruby[g]{質素}{しつそ}なり。朝廷\jruby[g]{賜}{たま}ふ所の\jruby[g]{賞典}{しやうてん}二千石は、\jruby[g]{悉}{こと{〴〵}}く私學校の\jruby[g]{費}{ひ}に\jruby[g]{充}{あ}つ。\jruby[g]{貧困}{ひんこん}なる者あれば、\jruby[g]{嚢}{のう}を\jruby[g]{傾}{かたぶ}けて之を\jruby[g]{賑}{すく}ふ。其の自ら視ること\jruby[g]{{※然}}{かんぜん}{\small ［＃※は「陷のつくり＋欠」、65-6］}として、\jruby[g]{微賤}{びせん}の時の如し。

\end{azconvIndented}

\begin{azconvIndented}{1\zw}{-1\zw}

九七　雅事多是虚、勿\azconvKunten{下}謂\azconvKunten{二}之雅\azconvKunten{一}而耽\azconvKunten{\azconvJR}之。俗事却是實、勿\azconvKunten{下}謂\azconvKunten{二}之俗\azconvKunten{一}而忽\azconvKunten{\azconvJR}之。

\end{azconvIndented}
\begin{azconvIndented}{2\zw}{-1\zw}

〔譯〕\jruby[g]{雅事}{がじ}多くは是れ\jruby[g]{虚}{きよ}なり、之を\jruby[g]{雅}{が}と謂うて之に\jruby[g]{耽}{ふけ}ること勿れ。俗事却て是れ實なり、之を俗と謂うて之を\jruby[g]{忽}{ゆるがせ}にすること勿れ。

\end{azconvIndented}

\begin{azconvIndented}{1\zw}{-1\zw}

九八　歴代帝王、除\azconvKunten{二}唐虞\azconvKunten{一}外、無\azconvKunten{二}眞禪讓\azconvKunten{一}。商周已下、秦漢至\azconvKunten{二}於今\azconvKunten{一}、凡二十二史、皆以\azconvKunten{レ}武開\azconvKunten{レ}國、以\azconvKunten{レ}文治\azconvKunten{レ}之。因知、武猶\azconvKunten{レ}質、文則其毛彩、虎豹犬羊之所\azconvKunten{二}以分\azconvKunten{一}也。今之文士、其可\azconvKunten{レ}忘\azconvKunten{二}武事\azconvKunten{一}乎。

\end{azconvIndented}
\begin{azconvIndented}{2\zw}{-1\zw}

〔譯〕\jruby[g]{歴代}{れきだい}の帝王、\jruby[g]{唐虞}{たうぐ}を\jruby[g]{除}{のぞ}く外、眞の\jruby[g]{禪讓}{ぜんじやう}なし。\jruby[g]{商周}{しやうしう}\jruby[g]{已下}{いか}\jruby[g]{秦漢}{しんかん}より今に至るまで、凡そ二十二史、皆武を以て國を開き、文を以て之を治む。因つて知る、武は猶\jruby[g]{質}{しつ}のごとく、文は則ち其の\jruby[g]{毛彩}{まうさい}にして、\jruby[g]{虎豹}{こへう}犬羊の分るゝ所以なるを。今の文士、其れ武事を忘る可けんや。

\end{azconvIndented}

\begin{azconvIndented}{1\zw}{-1\zw}

九九　遠方試\azconvKunten{レ}歩者、往往舍\azconvKunten{二}正路\azconvKunten{一}、※{\small ［＃※は「走にょう＋多」、66-3］}\azconvKunten{二}捷徑\azconvKunten{一}、或繆入\azconvKunten{二}林※{\small ［＃※は「くさかんむり／奔」、66-3］}\azconvKunten{一}、可\azconvKunten{レ}嗤也。人事多類\azconvKunten{レ}此。特記\azconvKunten{レ}之。

\end{azconvIndented}
\begin{azconvIndented}{2\zw}{-1\zw}

〔譯〕\jruby[g]{遠方}{えんぱう}に歩を\jruby[g]{試}{こゝろ}むる者、往往にして\jruby[g]{正路}{せいろ}を\jruby[g]{舍}{すて}て、\jruby[g]{捷徑}{せうけい}に\jruby[g]{{※}}{はし}{\small ［＃※は「走にょう＋多」、66-5］}り、或は\jruby[g]{繆}{あやま}つて\jruby[g]{{林※}}{りんまう}{\small ［＃※は「くさかんむり／奔」、66-5］}に入る、\jruby[g]{嗤}{わら}ふ可きなり。人事多く此に\jruby[g]{類}{るゐ}す。\jruby[g]{特}{とく}に之を\jruby[g]{記}{しる}す。

\end{azconvIndented}

\begin{azconvIndented}{1\zw}{-1\zw}

一〇〇　智仁勇、人皆謂\azconvKunten{二}大徳難\azconvKunten{\azconvIR}企。然凡爲\azconvKunten{二}邑宰\azconvKunten{一}者、固爲\azconvKunten{二}親民之職\azconvKunten{一}。其察\azconvKunten{二}奸慝\azconvKunten{一}、矜\azconvKunten{二}孤寡\azconvKunten{一}、折\azconvKunten{二}強梗\azconvKunten{一}、即是三徳實事。宜\azconvKunten{下}能就\azconvKunten{二}實迹\azconvKunten{一}以試\azconvKunten{レ}之可\azconvKunten{上}也。

\end{azconvIndented}
\begin{azconvIndented}{2\zw}{-1\zw}

譯智仁勇は、人皆\jruby[g]{大徳}{たいとく}\jruby[g]{企}{くはだ}て難しと謂ふ。然れども凡そ\jruby[g]{邑宰}{いふさい}たる者は、固と\jruby[g]{親民}{しんみん}の\jruby[g]{職}{しよく}たり。其の\jruby[g]{奸慝}{かんとく}を察し、\jruby[g]{孤寡}{こくわ}を\jruby[g]{矜}{あはれ}み、\jruby[g]{強梗}{きやうかう}を\jruby[g]{折}{くじ}くは、即ち是れ三徳の實事なり。宜しく能く實迹に就いて以て之を\jruby[g]{試}{こゝろ}みて可なるべし。

\end{azconvIndented}

\begin{azconvIndented}{1\zw}{-1\zw}

一〇一　身有\azconvKunten{二}老少\azconvKunten{一}、而心無\azconvKunten{二}老少\azconvKunten{一}。氣有\azconvKunten{二}老少\azconvKunten{一}、而理無\azconvKunten{二}老少\azconvKunten{一}。須\azconvKunten{丙}能執\azconvKunten{下}無\azconvKunten{二}老少\azconvKunten{一}之心\azconvKunten{上}、以體\azconvKunten{乙}無\azconvKunten{二}老少\azconvKunten{一}之理\azconvKunten{甲}。

\end{azconvIndented}
\begin{azconvIndented}{2\zw}{-1\zw}

譯身に\jruby[g]{老少}{らうせう}有りて、心に老少無し。氣に老少有りて、理に老少無し。須らく能く老少無きの心を\jruby[g]{執}{と}つて、以て老少無きの理を\jruby[g]{體}{たい}すべし。

\end{azconvIndented}
\begin{azconvIndented}{3\zw}{-1\zw}

〔評〕\jruby[g]{幕府}{ばくふ}南洲に\jruby[g]{禍}{わざはひ}せんと欲す。\jruby[g]{藩侯}{はんこう}之を\jruby[g]{患}{うれ}へ、南洲を\jruby[g]{大島}{おほしま}に\jruby[g]{竄}{ざん}す。南洲\jruby[g]{貶竄}{へんざん}せらるゝこと前後數年なり、而て身益\jruby[g]{壯}{さかん}に、氣益\jruby[g]{旺}{さかん}に、讀書是より大に進むと云ふ。

\end{azconvIndented}



\newpage
底本：「西郷南洲遺訓」岩波文庫、岩波書店

　　　1939（昭和14）年2月2日第1刷発行

　　　1985（昭和60）年2月20日第26刷発行

底本の親本：「南洲手抄言志録」博聞社

　　　1888（明治21）年5月17日発行

初出：「南洲手抄言志録」博聞社

　　　1888（明治21）年5月17日発行

※「「褒」の「保」に代えて「丑」」は「デザイン差」と見て「衰」で入力しました。

※底本の末尾に添えられた「書後の辭」で、秋月種樹氏が漢文の評言を附したとある。

入力：田中哲郎

校正：川山隆

2008年7月14日作成

2009年9月1日修正

青空文庫作成ファイル：

このファイルは、インターネットの図書館、青空文庫（http://www.aozora.gr.jp/）で作られました。入力、校正、制作にあたったのは、ボランティアの皆さんです。

\end{document}